\documentclass[10pt,a4paper]{amsart}
\usepackage[margin=19mm]{geometry}
\usepackage{amsmath,amssymb,mathtools}
\usepackage[scr=rsfs,cal=euler]{mathalfa}
\usepackage{xfrac}
\usepackage[unicode,hidelinks,bookmarks=false]{hyperref}
\newcommand{\ie}{i.e.\ }
\newcommand{\cf}{cf.\ }
\newcommand{\resp}{respectively\ }
\newcommand{\Subsection}{\subsection}
\providecommand{\nobreakdash}{\nobreak\hbox{-}}
\setlength{\emergencystretch}{2em}
\multlinegap0pt
\newcommand{\skpt}{\par\smallskip}
\providecommand{\Bar}{\overline}
\newcommand{\Changel}{\mathcode`l="7160}
\newcommand{\Changelback}{\mathcode`l="716C}
\newcommand{\NoChangel}[1]{%
\expandafter\let\csname old\string#1\endcsname=#1
\let#1=\relax
\newcommand{#1}{\mathop{\mathcode`l="716C\csname old\string#1\endcsname\mathcode`l="7160}\displaylimits}%
}
\newcommand{\NoChangelnolimits}[1]{%
\expandafter\let\csname old\string#1\endcsname=#1
\let#1=\relax
\newcommand{#1}{\mathop{\mathcode`l="716C\csname old\string#1\endcsname\mathcode`l="7160}\nolimits}%
}

\NoChangelnolimits{\log}
\NoChangelnolimits{\ln}
\NoChangel{\lim}
\NoChangel{\limsup}
\NoChangel{\liminf}
\NoChangel{\varlimsup}
\NoChangel{\varliminf}
\NoChangel{\varinjlim}
\NoChangel{\varprojlim}

\usepackage{tikz}
\usetikzlibrary{cd}
\tikzcdset{arrow style=Latin Modern}
\newcommand{\dradpl}{{\hspace*{-2mm}\begin{tikzcd}[ampersand replacement=\&,column sep = scriptsize]\mbox{}\ar[r,dashed]\&\mbox{}\end{tikzcd}\hspace*{-2mm}}}
\newcommand{\dratst}{{\hspace*{-2.5mm}\begin{tikzcd}[ampersand replacement=\&,column sep = small]\mbox{}\ar[r,dashed]\&\mbox{}\end{tikzcd}\hspace*{-2mm}}}
\let\drascst\dratst
\newcommand{\dra}{\mathchoice{\dradpl}{\dratst}{\drascst}{\drascst}}

\newenvironment{smallpmatrix}{\big(\begin{smallmatrix}}{\end{smallmatrix}\big)}

\newenvironment{Smallpmatrix}{\Big(\begin{smallmatrix}}{\end{smallmatrix}\Big)}

\newcommand\mto{\mathchoice{\longmapsto}{\mapsto}{\mapsto}{\mapsto}}

\AtBeginDocument{\renewcommand{\labelitemi}{$\scriptscriptstyle\bullet$}}

\newcommand{\Ab}{\textup{Ab}}
\newcommand{\id}{\textup{id}}
\newcommand{\Ind}{\textup{Ind}}
\newcommand{\op}{\textup{op}}
\newcommand{\quand}{\quad\text{and}\quad}
\newcommand{\sep}{\mkern2mu{:}\mkern2mu}
\DeclareMathOperator{\dett}{det}\let\det\dett
\DeclareMathOperator{\DET}{DET}
\DeclareMathOperator{\diag}{diag}
\DeclareMathOperator{\Gal}{Gal}
\DeclareMathOperator{\Imm}{Im}\let\Im\Imm
\DeclareMathOperator{\Mat}{Mat}
\DeclareMathOperator{\Nrd}{Nrd}
\DeclareMathOperator{\Sym}{Sym}

\let\epsilon\varepsilon

\usepackage[all,cmtip]{xy}\let\labelstyle\textstyle

\newtheorem{theorem}{Theorem}[section]
\newtheorem{lemma}[theorem]{Lemma}
\newtheorem{corollary}[theorem]{Corollary}
\newtheorem{proposition}[theorem]{Proposition}
\newtheorem{maintheorem}{Theorem}
\renewcommand\themaintheorem{\Alph{maintheorem}}

\theoremstyle{definition}
\newtheorem{definition}[theorem]{Definition}
\newtheorem{convention}[theorem]{Convention}
\newtheorem{remark}[theorem]{Remark}
\newtheorem{example}[theorem]{Example}
\newtheorem*{question}{Question}

\newcommand{\kk}{{\boldsymbol{k}}}
\newcommand{\FF}{{\boldsymbol{F}}}
\newcommand{\LL}{{\boldsymbol{L}}}

\DeclareMathOperator{\PGL}{PGL}
\DeclareMathOperator{\GL}{GL}
\DeclareMathOperator{\Bir}{Bir}
\DeclareMathOperator{\Aut}{Aut}
\DeclareMathOperator{\BirMori}{BirMori}


\title[Severi--Brauer groupoid determinant]{A determinant on birational maps of Severi--Brauer surfaces: original Corollary5.21}
\author{Elias Kurz}
\date{Source: J.\'E.P.--M.13(2026),1275--1322; DOI10.5802/jep.347}
\begin{document}\maketitle
\noindent\textit{Editorial scope.} The counted unit is the full original two-object groupoid Corollary5.21. The complete affine-descent, determinant-independence, trivial-relation, rewriting, six-link and splitting-field local core is retained as uncounted support. Borrowed original Figure1 on printed page1285 (physical PDF12) is an ancillary geometric overview and is not reproduced; every typed neighbor-point condition and six-link step remains. Native diagram mathematics and original generator, inverse and object conventions are unchanged. Original print slips, including the literal $\chi_i\in R$ and alternating automorphism types in the final proof, are retained without generated corrections.
\section{Original standing surface and groupoid setup}
Let $\kk$ be a perfect field. A~Severi--Brauer surface is a smooth projective surface $S$ over $\kk$ such that $S_{\bar{\kk}} \simeq \mathbb{P}^2_{\bar{\kk}}$. If $S$ is already isomorphic to $\mathbb{P}^2$ over $\kk$, we~say that $S$ is trivial. By a theorem of Châtelet this is equivalent to the existence of $\kk$-rational points. Let $\kk \subset \LL$ be a Galois extension such that $S$ splits (becomes trivial) over $\LL$. Choose an isomorphism $\varphi\colon S_{\LL} \simeq \mathbb{P}^2_{\LL}$. Then
\[
\varphi g \varphi^{-1} = \alpha_g g,\quad g \in \Gal(\LL/\kk), \;\alpha_g \in \PGL_3(\LL).
\]
Thus, under the isomorphism $\varphi$, the action of $g$ on $S_{\LL}$ corresponds to the action of $\alpha_g \circ g$ on $\mathbb{P}^2_{\LL}$. We~check that $g \mto \alpha_g$ is a cocycle. Moreover, sending the isomorphism class of $S$ to $\alpha_g$ gives the following one-to-one correspondence (\cite[I\S 5, III\S 1]{serre1997galois}):
\[
\left \{ \begin{aligned} & \text{ Severi--Brauer surfaces over~$\kk$ } \\[
-3pt] & \text{ that split over $\LL$ } \end{aligned} \right \} \longleftrightarrow H^1(\Gal(\LL/ \kk), \PGL_3(\LL)).
\]
Another classical correspondence is the following (which can be found, for example, in \cite[X\S 5]{serre2004corps}):
\[
\left \{ \begin{aligned} & \text{ Central simple algebras} \\[
-3pt] & \text{ of dimension 9 over~$\kk$ }\\[
-3pt]& \text{ that split over $\LL$ } \end{aligned} \right \} \longleftrightarrow H^1(\Gal(\LL/ \kk), \PGL_3(\LL)).
\]

Let $\LL$ be a splitting field of $S$ and let $\varphi$ be an isomorphism $\varphi\colon S_{\LL} \to \mathbb{P}^2_{\LL}$. For every $\alpha \in \Aut_{\kk}(S)$, we~have $ \varphi \alpha \varphi^{-1} \in \PGL_3(\LL)$. Up to multiplication by a scalar matrix, we~choose a representative $A \in \GL_3(\LL)$ of this automorphism such that (see Lemma~\ref{AutRep})
\[
A^g = A_g^{-1}AA_g, \quad g \in \Gal(\LL/\kk)
\]
for $A_g \in \GL_3(\LL)$ such that $[A_g] = \alpha_g$ (\cf Lemma \ref{3pLem}) and $A^g$ the image of the matrix $A$ under the action of $g$. This implies that $\det(A) \in \kk^*$ and we prove that its class in $\kk^*/(\kk^*)^3$ depends only on $\alpha$ and not on the representatives $A_g$ of $\alpha_g$, the isomorphism $\varphi$ or any splitting field (Lemma \ref{Determinant}). This construction therefore gives a group homomorphism
\[
\det \colon \Aut_{\kk}(S) \to \kk^*/(\kk^*)^3,
\]
which is the abelianization of the automorphism group of $S$ (Lemma \href{https://jep.centre-mersenne.org/articles/10.5802/jep.347/}{6.2}).

To extend this homomorphism to $\Bir_{\kk}(S)$, we~use the Sarkisov program, which decomposes elements of $\Bir_{\kk}(S)$ into simple birational maps between projective surfaces not necessarily isomorphic to $S$ (a detailed description can be found in Section~\ref{SubMori}). The Sarkisov program also embeds $\Bir_{\kk}(S)$ into a groupoid $\BirMori(S)$ (see Theorem \ref{Isk} for a definition). For Severi--Brauer surfaces, one has
\[\textstyle
\BirMori(S) = \{ S_1 \dra S_2: S_i \simeq S \text{ or } S_i \simeq S^{\op} \}.
\]
The surface $S^{\op}$ can be constructed using central simple algebras. Let $A$ be the central simple algebra corresponding to $S$. We~define $A^{\op}$ to be the algebra with the same underlying set but opposite multiplication, \ie $a \ast_{\op} b := b \ast a$. Thus, $S^{\op}$ is the Severi--Brauer surface corresponding to $A^{\op}$.

Using the Sarkisov program, one proves $\BirMori(S)$ is generated by the Sarkisov links based at $3$ and $6$\nobreakdash-points and automorphisms of $S,S^{\op}$. Here, a $d$-point is a closed point of degree $d$ which can be viewed as a Galois orbit of $d$ geometric points over~$\bar{\kk}$. The Sarkisov link based at a $3$\nobreakdash-point or $6$\nobreakdash-point $p$ is obtained by blowing up $p$ and contracting an orbit of $3$ or, respectively, $6$ disjoint $(-1)$-curves which are the strict transforms of the three lines through two of the three components or the six conics through five of the six components:
\[
\xymatrix{ & Y_6 \ar[dr]^{3} \ar[dl]_{3}& & & Y_3 \ar[dr]^{6} \ar[dl]_{6} & \\ S \ar@{-->}[rr] & & S^{\op} & S \ar@{-->}[rr] & & S^{\op}}
\]
The relations between the Sarkisov links are as follows: the trivial relations $\chi' \alpha \chi \beta = \id$ or $\chi \chi' = \id$ for fitting Sarkisov links $\chi, \chi'$ and automorphisms $\alpha,\beta$ and the elementary relations given in Theorem \ref{MainA}. This theorem shows that the determinant behaves well with the relations and is then the key point for extending the determinant from $\Aut_{\kk}(S)$ to $\Bir_{\kk}(S)$.

We consider the following subgroupoid of $\BirMori(S)$
\[
G_S := \{ S_1 \dra S_2 \mid S_i \in \{ S, S^{\op}\}\}.
\]
$G_S$ differs from $\BirMori(S)$ only because we do not allow different models of $S,S^{\op}$ to appear. Any map in $\BirMori(S)$ that starts or ends in a different model can be obtained by an element in $G_S$ composed with some isomorphisms between different models of $S,S^{\op}$.

Two Sarkisov links are said to be equivalent if they differ only by multiplication with automorphisms (\cf Definition \ref{def: equivLinks}). Let $\mathcal{E}_d$, $d \in \{3,6\}$ be the set of equivalence classes of $d$-links. We~choose a set $R$ of representatives for each class of $3$ or $6$\nobreakdash-links. We~extend the determinant to $G_S$ in the following way:
\begin{align*}
\det_R\colon G_S &\to \kk^*/(\kk^*)^3
\\
R &\mto 1
\\
\alpha \in \Aut_{\kk}(S) &\mto \det(\alpha)
\\
\beta \in \Aut_{\kk}(S^{\op})&\mto \det(\beta)^{-1}.
\end{align*}

\subsubsection*{Acknowledgements}
I would like to thank my thesis supervisor, Jérémy Blanc, for introducing me to this problem and for countless discussions. I would also like to thank the authors of \cite{BSY} for allowing me continue the work they had already done on this problem. The results taken from their unpublished notes (\cite{BSYUnp}) are Lemmas \ref{AutRep}, \ref{Fixpoints}, and \ref{DetTrivRef}, and Theorem \ref{Determinant}. I would like to thank Gabriel Dill for giving me some examples of field extensions of degree $3$ and results on their norm maps. I am grateful to the referee for their thorough reading and the many insightful suggestions that significantly improved the exposition.

\section{Preliminaries}

\subsection{Mori fiber spaces and Sarkisov links} \label{SubMori}
To study birational maps, we~use the Sarkisov program. The Sarkisov program has been studied extensively under various hypotheses and is an established tool for studying birational maps; see, for example, \cite{BSY}. We~use \cite{Lamy_2020} and \cite{Iskovskikh} as references for the Sarkisov program \cite{Lamy_2020} and \cite{Iskovskikh}; we refer to them for further details. We~first introduce rank-$r$ fibrations. We~are particularly interested in Mori fiber spaces because much can be said about their birational geometry.

\begin{definition} \label{RankRFibration}
Let $X$ be a smooth surface and $B$ a point or a smooth curve, both defined over a perfect field $\kk$. A~surjective morphism $\pi\colon X \to B$ is called \emph{rank-$r$ fibration} for $r \geq 1$, if:
\begin{itemize}
\item $-K_X$ is relatively ample, \ie $C\!\cdot\! (-K_X) \!>\! 0$ for all curves $C$ of $X$ contracted by~$\pi$.
\item The fibers of $\pi$ are connected.
\item The relative Picard rank $\rho(X/B)$ is equal to $r$.
\end{itemize}
A rank-$1$ fibration is also called \emph{Mori fiber space}.
\end{definition}
In the case where $B$ is a point, this is the same definition as that of a smooth del Pezzo surface of Picard rank $r$. When $B$ is a curve, the map $\pi$ is a conic bundle. For Severi--Brauer surfaces, we~only need the case in which $B$ is a point, $B = \{P\}$ below.

\begin{example} Here are a few examples of rank-$r$ fibrations.
\begin{itemize}
\item Let $X$ be a Severi--Brauer surface over $\kk$. we obtain a Mori fiber space:
\[
\xymatrix@R=.5cm{ X \ar[d] \\ \{P \} }
\]

\item Blowing up $X$ at a $3$\nobreakdash-point $p$ gives a del Pezzo surface of degree $6$ (all $3$\nobreakdash-points of~$X$ are in general position \cite[Lem.\,2.8]{zbMATH07206468}) which we denote by $Y$. Since we blow up one orbit, the rank increases by $1$, and we obtain the following rank-$2$ fibration:
\[
\xymatrix@C=.7cm@R=.5cm{& Y \ar[dl]_3 \ar[dd] \\ X \ar[d] & \\ \{P \} \ar[r]^{\simeq} & \{P\} }
\]

\item Let $Z$ be the blow-up of two $3$\nobreakdash-points of $X$ such that the six geometric points are in general position. Then, as in the previous example, we~obtain a rank-$3$ fibration:
\[
\xymatrix@C=.7cm@R=.5cm{&&Z \ar[ddd] \ar[dl]_3\\& Y \ar[dl]_3 \ar[dd] & \\ X \ar[d] & & \\ \{P \} \ar[r]^{\simeq} & \{P\} \ar[r]^{\simeq} & \{P\} }
\]
\end{itemize}
\end{example}
We now give a more formal definition of the construction in the example, by introducing the notion of a rank $r'$ fibration $X' \to B'$ dominating a rank-$r$ fibration $X \to B$. This will be useful for developing the tools needed to study the birational geometry of Mori fiber spaces.

\begin{definition}
Let $ X \to B, X' \to B'$ be rank $r$ and $r'$ fibrations, respectively, such that there exists a birational morphism $X' \to X$ and there exists a morphism $B \to B'$ such that:
\[
\xymatrix{X' \ar[r] \ar[d] & B' \\ X \ar[r] & B \ar[u] }
\]
We then say $X' \to B'$ \emph{dominates} $X \to B$. If the maps $X' \to X$, $B\to B'$ are both isomorphisms, we~say that $X' \to B'$, $X \to B$ are \emph{equivalent}.
\end{definition}
Note that $r' \geq r$, with equality if and only if the two fibrations are equivalent. We~now introduce the piece associated with a rank-$r$ fibration $X_r \to B_r$. This polytope records all rank $r'$ fibrations $X_{r'} \to B_{r'}$ dominated by $X_r \to B_r$. We~use it to study generators and relations of birational maps between Mori fiber spaces.

\begin{definition} \label{RPiece}
The \emph{piece} of a rank-$r$ fibration $X \to B$ is the $(r-1)$ dimensional combinatorial polytope constructed as follows:

Each rank $d$ fibration dominated by $X \to B$ is a $(d-1)$-dimensional face. For each pair of faces $X_i \to B_i$, $i=1,2$ the face $X_1 \to B_1$ is contained in $X_2 \to B_2$ if and only if the second fibration dominates the first one.
\end{definition}
The piece of a rank-$2$ fibration produces a birational map between two Mori fiber spaces. Those maps are called Sarkisov links and are the building blocks of birational maps between Mori fiber spaces.

\begin{definition} \label{SarkisovLink}
Let $X_2 \to B_2$ be a rank-$2$ fibration. By the two-rays game there are exactly two Mori fiber spaces $X \to B, X' \to B'$ dominated by $X_2 \to B_2$.
\[
\xymatrix@=.5cm{ & X_2 \ar[dl] \ar[dr] &\\X \ar@{-->}[rr] \ar[d]& & X' \ar[d] \\ B \ar[dr] & & B' \ar[dl] \\ & B_2 &}
\]
The induced map between $X$ and $X'$ is called \emph{Sarkisov link}. This is the piece of the fibration $X_2 \to B_2$.
\end{definition}
\begin{example}
Let $S$ be a Severi--Brauer surface and $p$ a $3$\nobreakdash-point. Let $L$ be the orbit of the $3$ (geometric) lines between any two of the three components of $p$. The map obtained by blowing up $p$ and contracting the strict transforms of the components of~$L$ is a Sarkisov link.

For a $6$\nobreakdash-point $q$ of $S$ take the (geometric) conics through any $5$ out of the $6$ components of $q$. They form an orbit, which we denote by $C$. The map obtained by blowing up $q$ and contracting the strict transforms of the components of $C$ is a Sarkisov link.
\[
\xymatrix{ & Y_6 \ar[dr] \ar[dl]& & & Y_3 \ar[dr] \ar[dl] & \\ S \ar@{-->}[rr] & & S^{\op} & S \ar@{-->}[rr] & & S^{\op}}
\]
\end{example}
The piece of a rank-$3$ fibration will induce a relation between Sarkisov links, which will play an essential role in determining the homomorphism we are looking for. An~example of such a relation of a Severi--Brauer surface can be found in Figure \href{https://jep.centre-mersenne.org/articles/10.5802/jep.347/}{1 in the original publication, p.1285}.

\begin{theorem}[{\cite[Prop.\,2.6]{Lamy_2020}}]\label{ElementaryRelation}
A $2$-piece is homeomorphic to a disk, and its boundary is a sequence of Sarkisov links, whose product is an automorphism.
\end{theorem}

We~say that this product encodes an \emph{elementary relation} between Sarkisov links. We~will use the following final result to study the group of birational self-maps of Severi--Brauer surfaces; it is known as the Sarkisov program.

\begin{theorem}[{\cite{Iskovskikh} and \cite[Th.\,3.1]{Lamy_2020}}] \label{Isk}
Let $X$ be a smooth projective surface over a perfect field $\kk$ that is birational to a Mori fiber space. The groupoid
\[
\BirMori(X) := \{X_1 \dra X_2 \text{ birational} \mid X_1,X_2 \text{ Mori fiber spaces birat. to } X\}
\]
is generated by Sarkisov links and isomorphisms of Mori fiber spaces. Any relation between Sarkisov links is generated by trivial and elementary relations.
\end{theorem}

The trivial relations are $\chi_1 \circ \chi_2 = \id$, where $\chi_i$ are Sarkisov links, inverse to each other, and $\alpha \circ \chi \circ \beta = \chi'$, where $\chi, \chi'$ are Sarkisov links and $\alpha, \beta $ are automorphisms of Mori fiber spaces.
\subsection{Properties of Severi--Brauer surfaces}
We consider Severi--Brauer surfaces, which are $\kk$-forms of $\mathbb{P}^2$. Their birational geometry has been studied (among others) by \cite{BSY}, \cite{zbMATH07206468}, \cite{cm:9040}.

\begin{definition}
Let $\kk$ be a perfect field. A~variety $S$ over $\kk$ is called a \emph{Severi--Brauer} variety if it is isomorphic to $\mathbb{P}^n$ over $\bar{\kk}$ for some $n \geq 0$. A~Severi--Brauer variety is called \emph{trivial} if it is isomorphic to $\mathbb{P}^n$ over $\kk$.
\end{definition}
By a theorem of Châtelet, a Severi--Brauer surface is trivial if and only if it has a $\kk$-rational point.
The first step in constructing our group homomorphism is to study the Sarkisov links on Severi--Brauer surfaces, as in \cite{BSY}. Another important tool in the study of Severi--Brauer surfaces is the orbit of points under the Galois action, also referred to as $d$-points.

\begin{lemma}[{\cite[Cor.\,2.2.2 with $p=3$]{BSY}}, {\cite{kollár2025severibrauervarietiesgeometrictreatment}}]
Let $S$ be a non-trivial Severi--Brauer surface over a perfect field $\kk$. Then $S$ does not contain points of degree $d$, where $d$ is not divisible by $3$. Moreover any non-trivial Severi--Brauer surface contains a point of degree $3$.
\end{lemma}

\begin{definition}
Let $p$ be a $3$\nobreakdash-point of $S$. Then the \emph{splitting field of $p$} is the smallest field $\LL \supseteq \kk$ such that the components of $p$ ($p = \{ p_1, p_2, p_3 \}$) are defined over~$\LL$. A~splitting field of the variety $S$ is any Galois extension $\LL$ of $\kk$ such that $S$ has $\LL$-rational points, $S(\LL) \neq \emptyset$.
\end{definition}

Let $\LL$ be the splitting field of a $3$\nobreakdash-point and $h \in \Gal(\bar{\kk}/\kk)$. Since $h(p) = p$, $p$~is defined over $h(\LL)$, which implies that $\LL \subset h(\LL)$ and by degree considerations $\LL = h(\LL)$. Thus, $\LL$ is a Galois extension of $\kk$. Since $S_{\LL}$ has $\LL$-rational points, $S_{\LL} \simeq \mathbb{P}^2_{\LL}$ and $\LL$ is also a splitting field of $S$.

\begin{lemma} [{\cite[Lem.\,2.3.2]{BSY}}] \label{3pLem}
Let $S$ be a Severi--Brauer surface over a perfect field $\kk$, $p$ a $3$\nobreakdash-point of $S$ and $\LL$ its splitting field. Then the following holds:
\begin{enumerate}
\item
The Galois group $\Gal(\LL/\kk)$ acts faithfully and transitively on $\{ p_1, p_2, p_3 \}$. In~particular, it is isomorphic to $\mathbb{Z}/3\mathbb{Z}$ or $\Sym_3$ and contains a unique $g \in \Gal(\LL/\kk)$ of order $3$ such that $g(p_1) = p_2, g(p_2) = p_3, g(p_3) = p_1$.
\item
There exists an $\LL$-isomorphism $\varphi \colon S_{\LL} \to \mathbb{P}^2_{\LL}$ such that\[
\varphi(p_1) = [1\sep 0\sep 0],\;\; \varphi(p_2) = [0\sep 1\sep 0],\;\; \varphi(p_3) = [0\sep 0\sep 1]\;\;\text{and}\;\;\varphi \circ g \circ \varphi^{-1} = A_g \circ g,
\]
with
$A_g = \begin{Smallpmatrix}
0 & 0 & \xi\\
1 & 0 & 0 \\
0 & 1 & 0
\end{Smallpmatrix}$
for some $\xi \in \LL^g$.
\item
If $q$ is a $3$\nobreakdash-point of $S$ having the same splitting field as $p$, there exists a $\alpha \in \Aut_{\kk}(S)$ that sends $p$ to $q$.
\end{enumerate}
\end{lemma}

As $-\mathrm{K}_S$ is ample and $\rho(S) = 1$, the surface $S$ is a Mori fiber space. We~now turn to the Sarkisov links of $S$. Every Severi--Brauer surface corresponds to a central simple algebra $A$ ({\cite[\S 6]{jahnel}}). We~define $A^{\op}$ to be the ring with the same set but reversed multiplication ($a \times_A b = b \times_{A^{\op}} a$). This algebra corresponds to another Severi--Brauer surface $S^{\op}$, which is not $\kk$-isomorphic to $S$ unless $S$ is trivial (see \cite[3.2.5/3.2.6]{BSY}). $S^{\op}$ is especially important because all Sarkisov links that start at $S$ end at $S^{\op}$.

\begin{lemma}[ {\cite[3.2.3, 3.2.7]{BSY}}, see also \cite{cm:9040}, \cite{zbMATH07206468}] \label{Links36}
Let $S$ be a non-trivial Severi--Brauer surface over a perfect field $\kk$ and let $\chi: S \dra S'$ be a Sarkisov link. The following are true:
\begin{enumerate}
\item The Sarkisov links $\chi, \chi^{-1}$ are centered on $d$-points, where $d \in \{3,6\}$. The two $d$-points at the center of $\chi,\chi^{-1}$ have the same splitting field.
\item We have $S' \simeq S^{\op}$ and any Mori fiber space birational to $S$ is isomorphic to either $S$ or $S^{\op}$.
\end{enumerate}

\end{lemma}
Composing an isomorphism with a Sarkisov link does not substantially change the nature of the link. We~therefore define equivalence classes of Sarkisov links, which are useful for constructing quotients of the group of birational self-maps.

\begin{definition} \label{def: equivLinks}
Let $\chi \colon X_1 \dra X_2$ and $\chi'\colon X'_1 \dra X'_2$ be two Sarkisov links between two rank $1$ del Pezzo surfaces over $\kk$. We~say $\chi, \chi'$ (or, respectively, their base points) are \emph{equivalent} if there exist isomorphisms $\alpha \colon X_1 \to X'_1$, $\beta \colon X_2 \to X'_2$ such that $\beta \circ \chi = \chi' \circ \alpha$.
\end{definition}

These equivalence classes are needed when working with the trivial relation \hbox{$\alpha \circ \chi \circ \beta = \chi'$} and are a key component in constructing the group homomorphism found in \cite{BSY}.

\begin{lemma} [{\cite[Lem.\,3.3.2]{BSY}}]
In the notation of Definition \ref{def: equivLinks}, one has the following:
\begin{enumerate}
\item\label{lem:2151}
$\chi, \chi'$ are equivalent.

\item\label{lem:2152}
There exists an isomorphism $\alpha$ as in Definition \ref{def: equivLinks} that sends the base points of~$\chi$ to the ones of $\chi'$.

\item\label{lem:2153}
There exists an isomorphism $\beta$ as in Definition \ref{def: equivLinks} that sends the base points of~$\chi^{-1}$ to the ones of $\chi'^{-1}$.
\end{enumerate}
\end{lemma}

\skpt
\begin{proof}
\eqref{lem:2151} $\Rightarrow$ \eqref{lem:2152}: Let $p,p'$ be the base points of $\chi,\chi'$. Since $\beta \circ \chi = \chi' \circ \alpha$ we obtain
\[
p=\Ind(\chi) = \Ind(\chi' \circ \alpha) = \alpha^{-1}(\Ind(\chi')) = \alpha^{-1}(p').
\]

\eqref{lem:2152} $\Rightarrow$ \eqref{lem:2153}: Set $\beta := \chi' \circ \alpha \circ \chi^{-1}$. We~claim that this is an automorphism; it then acts on the base points as described in \eqref{lem:2153}. We~write $\chi = \pi \circ \rho^{-1}$, $\chi' = \pi' \circ \rho'^{-1}$ for the minimal resolutions. Then since $\alpha$ sends the points blown-up in $\rho$ to the ones blown-up in $\rho'$, it follows that $\rho'^{-1} \circ \alpha \circ \rho$ is an automorphism sending the points blown-up in $\pi$ to the ones blown-up in $\pi'$, proving that $\beta$ is an automorphism.

\eqref{lem:2153} $\Rightarrow$ \eqref{lem:2151}: As in the second implication, we~prove $\alpha := \chi'^{-1} \circ \beta \circ \chi$ is an automorphism. This proves \eqref{lem:2151}.
\end{proof}

We define $\mathcal{E}_3, \mathcal{E}_6$ to be the sets of equivalence classes of $3$\nobreakdash-points and respectively $6$\nobreakdash-points or equivalently of $3$\nobreakdash-links and respectively $6$\nobreakdash-links. The next lemma gives the elementary relation between the Sarkisov links.

\begin{lemma}[{\cite[Lem.\,3.3.4]{BSY}}] \label{ElRel}
The elementary relations between Sarkisov links of a non-trivial Severi--Brauer surface $S$ defined over a perfect field $\kk$ are as follows:
\[
\chi_6 \circ \chi_5 \circ \chi_4 \circ \chi_3 \circ \chi_2 \circ \chi_1 = \id,
\]
where, for $i=1,2,3$, $\chi_{2i -1}\colon S \dra S^{\op}$ are links from the same class and $\chi_{2i}\colon S^{\op} \dra S$ are links from the same class.
\end{lemma}


Figure1 in the original publication (p.1285; not reproduced here) shows that each link maps the base points of the neighboring links to each other. This observation will be useful in the construction of the relation. The relation is given by the following hexagon (from \cite{BSY}):

\section{The determinant map}
To construct a useful group homomorphism that does not send the automorphisms to the neutral element, we~study how the automorphisms of a Severi--Brauer surface can be represented in the projective plane over a splitting field of $S$. We~will mainly use the splitting field of one fixed $3$\nobreakdash-point to obtain this representation. In~this chapter, we~study properties of the links and a useful representation of the birational maps of~$S$ in affine coordinates, which facilitate the study of the elementary relations and the automorphisms involved in them.

\begin{convention}
From now on let $\kk$ be a perfect field and $S$ a non-trivial Severi--Brauer surface over $\kk$.
\end{convention}

Conjugating an automorphism $\alpha$ of $S$ with an isomorphism to the projective plane over some splitting field gives an element of a projective linear group. These elements commute with the twisted action $A_g \circ g$, which imposes restrictions on the possible elements of the projective linear group. We~use these restrictions to find representative matrices with properties that are essential for constructing the group homomorphism on the automorphism group.

\begin{lemma}[\cite{BSYUnp}] \label{AutRep}
Let $Y$ be a Severi--Brauer variety defined over $\kk$ and $\kk \subset \LL$ be a finite Galois extension such that $Y_{\LL} \simeq \mathbb{P}^n_{\LL}$ via some isomorphism $\varphi$. There exist $[A_g] \in \PGL_{n + 1}(\LL)$, for all $ g \in \Gal(\LL/\kk)$, such that
\[
\varphi \circ g \circ \varphi^{-1} = [A_g] \circ g.
\]
Then for every $\alpha \in \Aut_{\kk}(Y)$ we can find a representative $A \in \GL_{n + 1}(\LL)$ of $\varphi \circ \alpha \circ \varphi^{-1}$ such that
\begin{equation} \label{eq:Rep}
A^g = A_g^{-1}AA_g,\quad \forall g \in \Gal(\LL/\kk).
\end{equation}
The matrix $A$ is unique up to $\kk^*$.
\end{lemma}

\begin{proof}
The $[A_g]$ exists since $\varphi \circ g \circ \varphi^{-1} \circ g^{-1}$ is an automorphism. Moreover, we~can observe that for $g,h \in \Gal(\LL/\kk)$ we have $[A_{gh}] = [A_g]g([A_h])$ and therefore there exists~$\mu_{gh}$ such that $A_{gh} = \mu_{gh}A_g g(A_h)$.

Let us choose a representative $A$ of $\varphi \circ \alpha \circ \varphi^{-1}$. Since $\alpha$ is defined over $\kk$, we find
\[
[A_g] \circ g \circ [A] = \varphi \circ g \circ \alpha \circ \varphi^{-1} = \varphi \circ \alpha \circ g \circ \varphi^{-1} = [A] \circ [A_g] \circ g,
\]
and thus $[A_g \circ A^g] = [A \circ A_g]$. This implies that there exists $\lambda_g \in \LL$ such that $\lambda_g A^g = A_g^{-1} A A_g$. We~calculate for $g,h \in \Gal(\LL/ \kk)$:
\[
A^{gh} \lambda_{gh} = A_{gh}^{-1} A A_{gh} = (A_h^{-1})^g A_g^{-1} A A_g A_h^g = \lambda_g (A_h^{-1} A A_h)^g = \lambda_g g(\lambda_h) A^{gh}.
\]
Thus $(g \mto \lambda_g) \in H^1(\Gal(\LL/\kk),\LL^*)$, but this only consists of the class of $1$\nobreakdash-cobound\-aries by Hilbert's Theorem 90. Thus, there exists $\lambda \in \LL^*$ such that $\lambda_g = \sfrac{\lambda}{g(\lambda)}$, replacing $A$ by $\lambda^{-1} A$ the claim follows.
\end{proof}
This lemma motivates the following definition, which we later extend to the case of birational maps. It is also the basis on which we define our determinant, as in the following theorem.

\begin{definition}
Let $Y$ be a Severi--Brauer variety defined over $\kk$ and $\kk \subset \LL$ be a finite Galois extension such that $Y_{\LL} \simeq \mathbb{P}^2_{\LL}$ via some isomorphism $\varphi$. Then there exist $[A_g] \in \PGL_{n + 1}(\LL)$, for all $ g \in \Gal(\LL/\kk)$, such that $\varphi \circ g \circ \varphi^{-1} = [A_g] \circ g$. The representative $A$ of $\varphi \circ \alpha \circ \varphi^{-1}$ satisfying \eqref{eq:Rep} is called an \emph{affine representative}.
\end{definition}

\begin{theorem}[\cite{BSYUnp}] \label{Determinant}
Let $Y$ be a Severi--Brauer variety of dimension $n$ defined over $\kk$ and $\alpha \in \Aut_{\kk}(Y)$ be an automorphism of $Y$. Choose an affine representative~$A$ of~$\alpha$. Then the following statements hold:
\begin{enumerate}
\item $\det(A) \in \kk^*$.
\item Letting $\det(\alpha)$ be the image of $\det(A)$ in $\kk^* / (\kk^*)^{n + 1}$ we obtain the following group homomorphism:
\begin{align*}
\Aut_{\kk}(Y) &\to \kk^*/(\kk^*)^{n + 1}
\\
\alpha &\mto \det(\alpha).
\end{align*}
Moreover this is independent of the field extension, the isomorphism and the choice of $A_g$.
\end{enumerate}
\end{theorem}

\begin{proof}
Since $A^g = A_g^{-1} A A_g$ we find
\[
g(\det(A)) = \det(A),\quad \forall g \in \Gal(\LL/\kk),
\]
and therefore $\det(A) \in \kk^*$. Reducing by $(\kk^*)^{n + 1}$ ensures that scaling $A$ by $\kk^*$ does not change the value, and thus the determinant does not depend on the choice of an affine representative.

For $\alpha, \beta$ with affine representatives $A,B$ we conclude that
\[
(AB)^g = A^g B^g = A_g^{-1}A A_g A_g^{-1} B A_g = A_g^{-1} AB A_g,
\]
Therefore, $AB$ represents $\alpha \circ \beta$. Thus, this determinant is a homomorphism. It remains to prove that the determinant depends only on $\alpha$.
\begin{itemize}
\item First, the choice of $A_g$ is unique only up to a scalar, which cancels out in $A_g^{-1} A A_g$. Thus the choice of $A_g$ does not affect the determinant.
\item If we fix $\LL$, for two possible isomorphisms $\varphi, \varphi'$, we~note that $\varphi' = \gamma \circ \varphi$ for $\gamma \in \Aut_{\LL}(\mathbb{P}^n)$. Using a representative $C$ of $\gamma$, we~note that we can now choose $C A_g (C^{-1})^g$ instead of $A_g$ and $CAC^{-1}$ instead of $A$. Therefore, we~conclude that
\begin{align*}
(CAC^{-1})^g = C^g A^g (C^{-1})^g&= C^g A_g^{-1} A A_g (C^{-1})^g
\\
&= (C A_g (C^{-1})^g)^{-1} CAC^{-1} (CA_g (C^{-1})^g).
\end{align*}
Therefore, $CAC^{-1}$ is our affine representative and the choice of the isomorphism $\varphi$ does not affect the determinant.
\item For different field extensions, we~first show the independence for $\LL \subset \FF$, where we use $\varphi$ defined on $\LL$ and an extended version of $\varphi$ on $\FF$.

We take $A, A_g \in \GL_{n + 1}(\FF)$, $g \in \Gal(\FF/\kk)$, satisfying \eqref{eq:Rep}. The elements in $\Gal(\FF/\LL)$ are elements from $\Gal(\FF/\kk)$ that restrict to the identity on $\LL$. Hence, for $g \in \Gal(\FF/\LL)$ we obtain
\[
[A_g] = \varphi \circ (\varphi^{-1})^g = \varphi \circ \varphi^{-1} = \id.
\]
therefore, we~choose $A_g = I_n$. We~then find $A^g = A_g^{-1} A A_g = A$, which proves that $A$ is defined over $\LL$.

Since every element $g \!\in\! \Gal(\LL/\kk)$ is the restriction of an element $g' \in \Gal(\FF/\kk)$ we get
\[
[A_g] = \varphi \circ (\varphi^{-1})^g = \varphi \circ (\varphi^{-1})^{g'} = [A_{g'}].
\]
Thus $A^g = A^{g'} = A_g^{-1} A A_g$, when we choose the same $A_g = A_{g'}$. Therefore, over $\LL$ we can choose the same affine representative and the determinant does not change.

\item Finally, for two extensions $\kk \subset \LL, \kk \subset \LL'$ using isomorphisms $\varphi, \varphi'$ we take $\FF:= \LL\LL'$ and use the independence for $\LL, \LL\LL'$ and $\LL', \LL\LL'$ and one time the independence for the extensions of $\varphi, \varphi'$ to $\LL\LL'$. \qedhere
\end{itemize}
\end{proof}

We want to define a map $\det \colon \Bir_{\kk}(S) \to \kk^* / (\kk^*)^3$ using the determinant. We~do this by using the Sarkisov program, \ie defining the determinant on Sarkisov links and automorphisms and proving that the map is well defined on the trivial and elementary relations.

\begin{definition}
Choose a representative $\chi_p$ of every equivalence class $p \in \mathcal{E}_3 \cup \mathcal{E}_6$ of $\kk$-links and collect them in the set $R$. Then $R,\Aut_{\kk}(S),\Aut_{\kk}(S^{\op})$ generates a subgroupoid of $\BirMori(S)$, which we call $G_S$. We~define:
\begin{enumerate}
\item for $\chi \in R$, $\det_R(\chi) := 1$;
\item for $\alpha \in \Aut_{\kk}(S)$, $\det_R(\alpha) := \det(\alpha)$;
\item for $\beta \in \Aut_{\kk}(S^{\op})$, $\det_R(\beta) := \det(\beta)^{-1}$.
\end{enumerate}
\end{definition}
We will fix the notation $R$ for the rest of the paper as the set of representatives of the equivalence classes of the Sarkisov links between $S$ and $S^{\op}$. We~want to extend this definition to a groupoid homomorphism
\[
\det_R\colon G_S \to \kk^*/(\kk^*)^3.
\]
To this end, we~use the fact that $G_S$ is generated by $R, \Aut_{\kk}(S), \Aut_{\kk}(S^{\op})$ with trivial and elementary relations. We~thus want to prove that for any trivial or elementary relation, the natural extension of this determinant map is sent to the identity. We~first prove that the trivial relation is sent to the identity, which will then imply that the determinant is well-defined on any Sarkisov link. Our trivial relations using only links in $R$ reduce to relations of the form $\chi = \alpha \chi \beta$, where $\chi$ is a Sarkisov link and $\alpha, \beta$ are automorphisms.

\begin{remark}
Let $\chi\colon S \dra S^{\op}$ be a Sarkisov link based on a $3$\nobreakdash-point. Over $\bar{\kk}$, $\chi$~is constructed by blowing-up the $3$\nobreakdash-point and contracting the three (geometric) lines through any two of the three components of the $3$\nobreakdash-point. This yields a map $\sigma$ (up to isomorphism) of the form
\begin{align*}
\sigma \colon \mathbb{P}^2_{\bar{\kk}}& \dra \mathbb{P}_{\bar{\kk}}^2\\
[x\sep y\sep z] &\mto [yz\sep  xz\sep  xy],
\end{align*}
which is called the standard Cremona involution.
\end{remark}

\begin{lemma} \label{BaseP}
Let $\chi\colon S \dra S^{\op}$ be a Sarkisov link based at $p$ and let $p^{\op}$ be the base point of $\chi^{-1}$. Then there exists a splitting field $\LL$ of $S,S^{\op}$ and isomorphisms
\[
\varphi\colon S_{\LL} \to \mathbb{P}^2_{\LL}, \quad \varphi^{\op}\colon S_{\LL}^{\op} \to \mathbb{P}^2_{\LL}
\]
such that the following statements hold:
\begin{enumerate}
\item \label{case: 1bp} If $p,p^{\op}$ are $3$\nobreakdash-points $\psi := \varphi^{\op} \circ \chi \circ \varphi^{-1} = D\sigma$, where $D$ is a diagonal matrix over $\LL$.
\item \label{case: 2bp} If $p,p^{\op}$ are $6$\nobreakdash-points, $\psi := \varphi^{\op} \circ \chi \circ \varphi^{-1}$ is the blow-up of $6$ points and the blow-down of the $6$ conics through $5$ of the $6$ points.
\end{enumerate}
\end{lemma}

\skpt
\begin{proof}
\begin{itemize}
\item If we blow-up a point of degree $3$ in $S$ we have a del Pezzo surface $Y$ of degree~$6$ which has six $(-1)$-curves. The three exceptional divisors which form an orbit and the strict transforms of the geometrical lines through two of the three components which form an orbit. Therefore, we~contract the strict transforms to obtain $\chi$. We~can choose $\varphi, \varphi^{\op}$ so that it maps $p,p^{\op}$ to the coordinate points. Thus this construction of $\chi$ yields $\psi = D \sigma$ for a diagonal matrix $D$ and the standard Cremona involution $\sigma$.
\item If we blow-up a point of degree $6$, we~obtain a degree $3$ del Pezzo surface $Y$ with $27$ $(-1)$-curves. Six of them are the orbit made up by the exceptional divisors, and another six are the strict transforms of the (geometric) conics through five of the six components of our $6$\nobreakdash-point $p$. The other $15$ $(-1)$ curves come from the strict transforms of the geometric lines between two of the six points. A~small combinatorial argument shows that they are not contractable because they do not have disjoint orbits. Therefore, the Sarkisov link based at the $6$\nobreakdash-point is obtained by contracting the strict transforms of the six (geometric) conics.
\end{itemize}
\end{proof}

In the case of a $3$\nobreakdash-point the splitting field $\LL$ and isomorphisms $\varphi, \varphi^{\op}$ are as given in Lemma \ref{3pLem}, where we use $\xi^{-1}$ as the parameter for $S^{\op}$. Next we look at the image of $3$\nobreakdash-points under Sarkisov links.

\begin{lemma} \label{ImageOf3P}
Let $p,q$ be two different $3$\nobreakdash-points of S and $\chi$ a Sarkisov link based at~$q$. Then $\chi(p)$ is a $3$\nobreakdash-point with the same splitting field as $p$.
\end{lemma}

\begin{proof}
Let $g \in G:= \Gal(\bar{\kk}/\kk)$, since $\chi$ is $g$-invariant we have $g(\chi(p_i)) = \chi(g(p_i))$. This equation implies that $\chi(p)$ is a $d$-point of orbit size smaller than $3$. Since the order of all points of $S^{\op}$ is divisible by $3$ we obtain that $\chi(p)$ is a $3$\nobreakdash-point.

Let $\LL$ be the splitting field of $p$ and let $\LL'$ be the one of $\chi(p)$. Since $p$ is defined over $\LL$ and $\chi$ is defined over $\kk$ we conclude that $\chi(p)$ is defined over $\LL$, thus $\LL' \subseteq \LL$. We~know that $\kk \subset\LL$ is of degree $3$ or $6$ with the Galois group $\Sym_3$. Since the same holds for $\kk \subset \LL'$ and we conclude that $\Sym_3$ does not have a quotient of order $3$ we also find that $\LL$ has no strict subfield that is a degree-$3$ Galois extension of $\kk$.
\end{proof}
A birational map of $\mathbb{P}^2$ is given by three homogeneous coprime polynomials of the same degree. The degree of the birational map is defined as the degree of these polynomials. Using Lemma \ref{BaseP} we can study the degree of birational maps between~$S$ and $S^{\op}$ over $\LL$.

\begin{remark} \label{rem: HomSys}
Since $S_{\bar{\kk}} \simeq \mathbb{P}^2_{\bar{\kk}}$ we introduce tools for studying birational maps of the projective plane. Let $\theta$ be any birational map $\theta \colon \mathbb{P}^2 \dra \mathbb{P}^2$. Its homaloidal system is the pullback of the pencil of lines under $\theta$. The characteristic of $\theta$ is given by $d;m_1, \dots , m_r$, where $d$ is the degree of the map $\theta$ and $m_i$ are the multiplicities of the base points of $\theta$ in its homaloidal system. For further details, see \cite[\S 2]{zbMATH01683518}. We~now have the following three equalities, which are called Noether's relations \cite[\S 2.5, p.\,51]{zbMATH01683518}.
\begin{enumerate}
\item $\sum_{i = 1}^r m_i^2 = d^2 - 1$,

\item $\sum_{i = 1}^r m_i = 3(d - 1)$,

\item $\sum_{i = 1}^r m_i (m_i - 1) = (d - 1) (d - 2)$.
\end{enumerate}
Take $\alpha_1, \alpha_2, \alpha_3$ to be the multiplicities in the homaloidal system of the following points:
\[
p_1 := [1\sep 0\sep 0],\quad p_2 := [0\sep 1\sep 0],\quad p_3 := [0\sep 0\sep 1].
\]
We can use \cite[\S 4]{zbMATH01683518} to calculate the characteristic of $\theta \circ \sigma$. We~find the multiplicities of the base points of $\theta$ remain the same for all points except $p_1,p_2,p_3$. The points $p_1,p_2,p_3$ have multiplicities
\[
d - \alpha_2 - \alpha_3,\quad d - \alpha_1 - \alpha_3,\quad d - \alpha_1 - \alpha_2.
\]
By Noether's second relation, we~observe that
\[
d = \deg(\theta) = \frac{1}{3} \Bigl(\sum_{i = 1}^r m_i\Bigr) + 1.
\]
Using this for $\theta \circ \sigma$ also shows that the difference in degree of the two birational maps is a third of the difference in multiplicities, \ie
\[
\deg (\theta \circ \sigma) = d + (d - \alpha_1 - \alpha_2 - \alpha_3) = 2d - \alpha_1 - \alpha_2 - \alpha_3.
\]
\end{remark}

\begin{lemma} \label{lem: degofbir}
Let $p$ be a $3$\nobreakdash-point of $S$ and $q$ be a $3$\nobreakdash-point of $S^{\op}$ with the same splitting field $\LL$. We~choose isomorphisms $\varphi \colon S_{\LL}\to \mathbb{P}^2_{\LL}$ and $\varphi^{\op}\colon S^{\op}_{\LL} \to \mathbb{P}^2_{\LL}$ for $p,q$ as in Lemma \ref{3pLem}.
Then for every $\psi \colon S \dra S^{\op}$ we have
\[
\deg(\varphi^{\op} \circ \psi \circ \varphi^{-1}) \equiv 2 \pmod 3,
\]
and, for $\psi \colon S \dra S$,
\[
\deg(\varphi \circ \psi \circ \varphi^{-1}) \equiv 1 \pmod 3.
\]
\end{lemma}

\begin{proof}
We start by noting that the result is true for automorphisms. We~can write $\psi = \chi_k \circ \dots \circ \chi_1$, where $\chi_i$ are Sarkisov links and assume that this decomposition is minimal. Furthermore, we~use the fact that all components have the same multiplicity. We~prove the statement by induction on the length $k$ of a minimal decomposition into Sarkisov links. Note that $\psi \circ \chi_1^{-1} = \chi_k \circ \dots \circ \chi_2$ is a minimal decomposition of length $k-1$.

\subsubsection*{$k=1$}
If $\psi = \chi_1$ we conclude that $\varphi^{\op} \circ \chi_1 \circ \varphi^{-1}$ is of characteristic $2;1^3$ or $5;2^6$, by Lemma \ref{BaseP}. We~find $2 \equiv 2 \pmod 3$ and $5 \equiv 2 \pmod 3$. We~also find that the components of base points belonging to one $3$\nobreakdash-point or $6$\nobreakdash-point have the same multiplicity.

\subsubsection*{$k \to k+1$}
We define $\sigma_i := \varphi^{\op} \circ \chi_i \circ \varphi^{-1}$ for $i$ odd and $\sigma_i := \varphi \circ\chi_i \circ (\varphi^{\op})^{-1}$ for $i$ even. Then we need to look at $\sigma_{k+1} \circ \dots \circ \sigma_1$ by induction $\deg(\sigma_{k+1} \circ \dots \circ \sigma_2) = 3r + \delta$, where $\delta = 2$ if $k$ is odd and $\delta = 1$ if $k$ is even.

If $\sigma_1$ is as in (\ref{case: 1bp}) we write $\alpha$ for the multiplicity of the base points of $\sigma_1^{-1}$ in $\sigma_{k+1} \circ \dots \circ \sigma_2$. Then by Remark \ref{rem: HomSys} we get
\[
\deg(\sigma_{k+1} \circ \dots \circ \sigma_1) = 2(3r + \delta) - 3\alpha = 3(2r - \alpha) + 2 \delta,
\]
which is equivalent to $1$ modulo $3$ if $\delta = 2$ and to $2$ modulo $3$ if $\delta = 1$. The multiplicities of the base points of $\sigma_1^{-1}$ are now $3r + \delta - 2 \alpha$, so they are equal. The multiplicities of the other points remain the same.

Now, if $\sigma_1$ is as in (\ref{case: 2bp}). Then $\sigma_1$ is the blow-up of $6$ points and the blow down of the $6$ conics through, respectively, $5$ of the points. Let $\alpha$ be the multiplicity of the base points of $\sigma_1^{-1}$ in $\sigma_{k+1} \circ \dots \circ \sigma_2$, then we conclude that the intersection of the strict transform of an element $C$ of the homaloidal map of $\sigma_{k+1} \circ \dots \circ \sigma_2$ with the strict transform of the conics is $2(3r + \delta) - 5\alpha$, which is the new multiplicity of those points.

The multiplicities of the other points remain unchanged. Using Noether's second relation we conclude that
\[
\deg(\sigma_{k+1} \circ \dots \circ \sigma_1) = 5(3r + \delta) - 12 \alpha = 3(5r - 4 \alpha) + 5\delta,
\]
which is equivalent to $1$ if $\delta = 2$ and equivalent to $2$ if $\delta = 1$. For a more detailed proof of this fact, we~refer to \cite[\S 4]{zbMATH01683518}.
\end{proof}
\begin{convention} \label{conv: phi}
For a splitting field $\LL$ of a $3$\nobreakdash-point $p := \{ p_1, p_2, p_3\}$ of $S$ we always take $g$ to be the element of the Galois group such that $g(p_1) = p_2$, $g(p_2) = p_3$, $g(p_3) = p_1$. We~write $\varphi$ for an $\LL$-isomorphism between $S_{\LL}$ and $\mathbb{P}^2_\LL$ with the properties given in Lemma \ref{3pLem}. We~also write $A_g, \xi$ for the respective objects defined in Lem\-ma~\ref{3pLem}.

By \cite[Lem.\,3.2.9]{BSY}, we~can choose for $S^{\op}$ a similar isomorphism denoted by~$\varphi^{\op}$ and a similar matrix denoted by $A_{g}^{\op}$ given by the parameter $\xi^{-1}$.
\end{convention}

\begin{proposition} \label{AffineRep}
Let $p$ be a $3$\nobreakdash-point of $S$ whose splitting field $\LL$ satisfies \hbox{$[\LL:\kk]=3$}. We~then take $\varphi, \varphi^{\op}$ as in Convention \ref{conv: phi}. We~may assume that $\Gal({\LL/\kk})$ is generated by our $g$.

For $\psi \colon S \dra S^{\op}$ we find $d_1 \geq 0$, $F_0,F_1,F_2 \in \LL[x,y,z]$ homogeneous of degree $3d_1 + 2$ such that $\varphi^{\op} \circ \psi \circ \varphi^{-1} = [F_0\sep F_1\sep F_2]$ and:
\begin{equation} \label{eq:RepB}
\xi^{i(d_1+1)} (F_0^{g^i}, F_1^{g^i}, F_2^{g^i}) = (A_{g^i}^{\op})^{-1} \circ (F_0,F_1,F_2) \circ A_{g^i}, \quad i=0,1,2.
\end{equation}
For $\psi \colon S \dra S$, we~find $d_2 \geq 0$, $G_0,G_1,G_2 \in \LL[x,y,z]$ homogeneous of degree $3d_2 + 1$ such that $\varphi \circ \psi \circ \varphi^{-1} = [G_0\sep G_1\sep G_2]$ and:
\begin{equation}
\label{eq:RepB2} \xi^{i\cdot d_2} (G_0^{g^i}, G_1^{g^i}, G_2^{g^i}) = A_{g^i}^{-1} \circ (G_0,G_1,G_2) \circ A_{g^i}, \quad i=0,1,2.
\end{equation}
\end{proposition}

\begin{proof}
By Lemma \ref{3pLem}, we~can take
\[\arraycolsep4.5pt
A_g =
\begin{pmatrix}
0 & 0 & \xi \\
1 & 0 & 0 \\
0 & 1 & 0
\end{pmatrix}\!.
\]
We then also take $A_{id} = I_3$, $A_{g^2} = A_g^2$, then $A_g^3 = \xi I_3$ and similar for $A_g^{\op}$ just with~$\xi^{-1}$ instead of $\xi$. We~will just do the first case, and the second will work similarly.

We take $F := \varphi^{\op} \circ \psi \circ \varphi^{-1} = [F_0\sep F_1\sep F_2]$, where $F_i \in \LL[x,y,z]$ are homogeneous, coprime of degree $3d_1 + 2$ (possible due to Lemma \ref{lem: degofbir}). Since $\psi$ is $\kk$-birational, we~get $F \circ A_{g^i} \circ g^i = A_{g^i}^{\op} \circ g^i \circ F$ which gives $F^{g^i} = g^{i} \circ F \circ g^{-i} = (A_{g^{i}}^{\op})^{-1} \circ F \circ A_{g^{i}}$. Therefore, there exists $\mu_{g^{i}} \in \LL^*$ such that
\[
(A_{g^{i}}^{\op})^{-1} \circ (F_0, F_1, F_2) \circ A_{g^{i}} = (F_0^{g^{i}},F_1^{g^{i}},F_2^{g^{i}})\mu_{g^{i}},
\]
where $\mu_{id} = 1$.

Next, we~show that $(g^{i} \mto \sfrac{\mu_{g^i}}{\xi^{i(d+1)}})$ defines a $1$-cocycle.
For $u_1,u_2$ in the Galois group there exist $\lambda_{u_1,u_2}, \lambda_{u_1,u_2}^{\op} \in \LL^*$ such that
\[
A_{u_1 \circ u_2} = \lambda_{u_1,u_2} A_{u_1} u_1(A_{u_2}), \quad A_{u_1 \circ u_2}^{\op} = \lambda_{u_1,u_2}^{\op} A_{u_1}^{\op} u_1(A_{u_2}^{\op})
\]
Since $\Gal(\LL/\kk) = \{id,g,g^2 \}$ our choice of $A_{g^i}$ implies that all $\lambda_{u_1,u_2} = 1$ except for
\[
\lambda_{g^2,g^2} = \lambda_{g^2,g} = \lambda_{g,g^2} = \xi^{-1}
\quand
\lambda_{g^2,g^2}^{\op} = \lambda_{g^2,g}^{\op} =\lambda_{g,g^2}^{\op} = \xi.
\]
We then obtain
\begin{align*}
\mu_{u_1u_2}(F_0^{u_1u_2},F_1^{u_1u_2},F_2^{u_1u_2}) &= (\lambda_{u_1,u_2}^{\op} A_{u_1}^{\op} u_1(A_{u_2}^{\op}))^{-1} {\circ} (F_0,F_1,F_2) {\circ} (\lambda_{u_1,u_2} A_{u_1} u_1(A_{u_2}))
\\&
= \frac{\lambda_{u_1,u_2}^{\deg(F)} \mu_{u_1}}{\lambda_{u_1,u_2}^{\op}} u_1((A_{u_2}^{\op})^{-1} {\circ} (F_0, F_1, F_2) {\circ} A_{u_2})
\\&
= \frac{\lambda_{u_1,u_2}^{\deg(F)} \mu_{u_1} u_1(\mu_{u_2})}{\lambda_{u_1,u_2}^{\op}} (F_0^{u_1u_2},F_1^{u_1u_2},F_2^{u_1u_2}).
\end{align*}
Thus,
\[
\frac{\lambda_{u_1,u_2}^{3d + 2} \mu_{u_1} u_1(\mu_{u_2})}{\lambda_{u_1,u_2}^{\op}} = \mu_{u_1,u_2}.
\]
Moreover, the fraction $\sfrac{\lambda_{u_1,u_2}^{3d + 2}}{\lambda_{u_1,u_2}^{\op}}$ is equal to $1$, except in the case where $(u_1,u_2) \in \{ (g,g^2), (g^2,g),(g^2,g^2) \}$, in which it is equal to $\xi^{-3(d+1)}$. We~thus obtain the following equations:
\begin{gather*}
\frac{\mu_{g^2}}{\xi^{2(d+1)}} = \frac{\mu_g}{\xi^{d+1}} \frac{g(\mu_g)}{\xi^{d+1}},
\\
\mu_{\id} = \frac{\mu_g}{\xi^{(d+1)}} \frac{g(\mu_{g^2})}{\xi^{2(d+1)}} = \frac{\mu_{g^2}}{\xi^{2(d+1)}} \frac{g^2(\mu_g)}{\xi^{d+1}}, \quad \frac{\mu_g}{\xi^{d+1}} = \frac{\mu_{g^2}}{\xi^{2(d+1)}}\frac{g^2(\mu_{g^2})}{\xi^{2(d+1)}}.
\end{gather*}
We observe that $(g^{i} \mto \sfrac{\mu_{g^i}}{\xi^{i(d+1)}}) \in H^1(\Gal(\LL/\kk),\LL^*)$ which is trivial by Hilbert's Theorem 90. Therefore, there exists $\epsilon \in \LL^*$ such that $\sfrac{\mu_{g^{i}}}{ \xi^{i(d+1)}} = \sfrac{\epsilon}{g^{i}(\epsilon)}$. Replacing~$F_i$ with $\epsilon^{-1} F_i$ gives the result.
\end{proof}
Using this proposition, we~can now generalize the definition of affine representatives to birational maps. For automorphisms, this agrees with the previous definition by identifying linear automorphisms of $\mathbb{A}^3$ with invertible matrices.

\begin{definition}
Let $p$ be a $3$\nobreakdash-point of $S$ such that for its splitting field $\LL$ we have $[\LL\sep \kk]=3$. Take the conventional $\varphi, \varphi^{\op}$ for $p$ and a $3$\nobreakdash-point $q$ of $S^{\op}$ with the same splitting field. We~may assume that $\Gal({\LL/\kk})$ is generated by our $g$. Let $\psi \colon S \dra S'$, $S' \in \{S,S^{\op}\}$ be a birational map. For the $F_i \in k[x,y,z]$ found in Proposition \ref{AffineRep} the map $\mathbb{A}^3 \to \mathbb{A}^3, x \mto (F_0(x),F_1(x),F_2(x))$ is called an \emph{affine representative}.
\end{definition}

\begin{example}
Let $\sigma$ be the standard Cremona involution that has coordinates \hbox{$[yz\sep xz\sep xy]$}. We~define $\Sigma := (yz, xz, xy) \in \LL[x,y,z]^3$. A~direct calculation shows that~$\Sigma$ satisfies the equation in Proposition \ref{AffineRep} and therefore $\Sigma$ is an affine representative of a birational map from $S$ to $S^{\op}$. We~denote this affine representative by~$\Sigma$ throughout.
\end{example}
The affine representatives give a way to describe all automorphisms of $S$ defined over $\kk$. Moreover, their affine representatives depend on exactly three parameters, as~proved by the following result.

\begin{lemma} \label{matrixL}
The affine representatives of the automorphisms of $S$ correspond to the following matrices:
\[\arraycolsep4.5pt
\begin{pmatrix}
a & \xi g(c) & \xi g^2(b) \\
b & g(a) & \xi g^2(c) \\
c & g(b) & g^2(a)
\end{pmatrix}\!,
\]
where $(a,b,c) \in \LL^3 \setminus \{ (0,0,0) \}$.
\end{lemma}

\begin{proof}
Let $\alpha \in \Aut_{\kk}(S)$. Proposition \ref{AffineRep} gives the same result as Lemma \ref{AutRep}, thus that the affine representative $A$ of $\alpha$ satisfies $A^g = A_g^{-1} A A_g$. Let $a_i$ be the columns of $A$. We~can observe that
\begin{align*}
A^g = A_g^{-1} A A_g &\iff A_g(g(a_1)) = a_2, A_g(g(a_2)) = a_3, A_g(g(a_3)) = \xi a_1
\\
&\iff a_2 = A_g(g(a_1)), a_3 = A_g^2(g^2(a_1)).
\end{align*}
Thus our matrix is given by:
\begin{equation*}\arraycolsep4.5pt
\begin{pmatrix}
a & \xi g(c) & \xi g^2(b) \\
b & g(a) & \xi g^2(c) \\
c & g(b) & g^2(a)
\end{pmatrix}\!,
\end{equation*}
where $a_1 = (a,b,c) \neq (0,0,0)$. This matrix is always invertible and thus always an affine representative, since for any $a_1 \neq (0,0,0)$ the $A_g \circ g$ orbit of $a_1$ corresponds to a $3$\nobreakdash-point of $S$ which is always in general position by Lemma \ref{3pLem}.
\end{proof}
Every $\LL$-birational map of $\mathbb{P}^2$ that has such affine representatives corresponds via~$\varphi$ to a $\kk$-birational map of $S$, since it does not change under the Galois action. We~also note that such affine representatives always come with $\LL,\varphi, \varphi^{\op}, A_g$. This only works in splitting fields that are degree-$3$ extensions, and thus we show in the next lemma that such a field always exists.

\begin{lemma} \label{lem: goodSplit}
There exists a $3$\nobreakdash-point of $S$ whose splitting field $\LL$ is a degree-$3$ extension of $\kk$.
\end{lemma}

\begin{proof}
With Lemma 2.3.4 in \cite{BSY} we find a field $\LL$ such that the field extension of $\kk \subseteq \LL$ is a degree-$3$ Galois extension and $S$ has points over $\LL$ (it is isomorphic to~$\mathbb{P}^2$ over $\LL$). We~want to prove that $\LL$ is a splitting field of a $3$\nobreakdash-point.

We take the preimage of $[1\sep 0\sep 0]$ with respect to the isomorphism. Then we take the Galois orbit of this point with respect to $\Gal(\LL/\kk)$. This is then a $d$-point for $d \leq 3$. Therefore, it must be a $3$\nobreakdash-point, since there are no $2$-points or $1$-points in $S$. We~then get that the splitting field is in $\LL$ and is not $\kk$, thus it must be $\LL$, since $3$ is prime.
\end{proof}

Given an automorphism of $S$, we aim to obtain affine representatives that are easy to handle. To~this end, we~prove that they always fix three geometric points that form a $3$\nobreakdash-point, and therefore we can choose an $\LL$-isomorphism $\varphi \colon S_{\LL} \to \mathbb{P}^2_{\LL}$ such that the affine representative becomes a diagonal matrix.

\begin{lemma}[\cite{BSYUnp}] \label{Fixpoints}
Let $\alpha \in \Aut_{\kk}(S) \setminus \{ \id \}$. Then there exists exactly one $3$\nobreakdash-point in $S$ whose components are fixed by $\alpha$ and its components are the only $\bar{\kk}$-points, fixed by $\alpha$.
\end{lemma}

\begin{proof}
We take the degree-$3$ splitting field $\LL$ of a point $p$ from Lemma \ref{lem: goodSplit} and the according $\varphi$ from Lemma $\ref{3pLem}$. We~now choose $A$ to be an affine representative of $\alpha$ and see that its eigenvectors correspond exactly to the fixed points of $\alpha$. Those do not necessarily need to be defined over $\LL$, but since the eigenvalues come from a degree-$3$ polynomial, we~can take $\LL'$ a Galois extension of $\kk$ containing $\LL$ such that they are defined. We~extend $\varphi$ to $\LL'$ and replace $A$ by the affine representative of $\alpha$ over $\LL'$ (possibly with a different scalar due to the new Galois group).

All eigenvalues of $A$ are defined over $\LL'$ and $A$ has an eigenvector defined over~$\LL'$ for any eigenvalue. Therefore, take $v, \lambda$ an eigenvector and its eigenvalue. We~get for every $u \in \Gal(\LL'/ \kk)$ $A_u A^u = A A_u$ and thus $A_u u(v)$ is an eigenvector of $A$ with eigenvalue $u(\lambda)$.

Take $V_1, \dots , V_r$ the eigenspaces of $A$ over $\LL'$. And take $V$ the vector space generated by all of them. We~then have $A_u u(V) = V$. Therefore, we~conclude that the preimage of the projectivization of $V$ under $\varphi$ is a twisted linear subvariety, and thus $V$ has dimension $3$.

There are three cases:
\begin{enumerate}
\item $r = 1$ and therefore $V = V_1$, which is only possible for the identity.
\item If we have $r=2$, then by a reordering we may assume that $\dim(V_1) =\nobreak 2$, $\dim(V_2) = 1$. For $v, w$ a basis of $V_1$ with eigenvalues $\lambda_1$, $A_u u(v), A_u u(w)$
are eigenvectors of eigenvalue $u(\lambda)$. But $\dim(V_2) = 1$, so they are in $V_1$ and therefore $A_uu(V_1)) =\nobreak V_1$. This makes the preimage of the projectivization of $V_1$ a twisted linear subvariety of dimension $1$, which is impossible.
\item The only remaining option is $r=3$ and $\dim(V_i) = 1$. But then we conclude that the preimages of the projectivization of $V_i$ are three $\LL'$-points, which are Galois invariant as a set, and therefore must form a $3$\nobreakdash-point. Moreover, any eigenvector of $A$ over $\bar{\kk}$ is a multiple of an eigenvector defined over $\LL'$ and therefore their projectivizations have the same class.\qedhere
\end{enumerate}
\end{proof}

We now have all the tools to study how the trivial relations behave under our generalization of the determinant map.

\begin{lemma}[\cite{BSYUnp}] \label{DetTrivRef}
If for $\alpha \in \Aut_{\kk}(S)$, $\beta \in \Aut_{\kk}(S^{\op})$ and $\chi$ a Sarkisov link of~$S$ we have $\beta \chi = \chi \alpha$ then $\det(\beta) = \det(\alpha)^{-1}$.
\end{lemma}

\begin{proof}
If we choose a different representative $\chi' = \gamma_1 \chi \gamma_2$ for $\kk$-automorphisms $\gamma_1, \gamma_2$ of $S^{\op},S$, we~can replace $\beta, \alpha$ by $\gamma_1 \beta \gamma_1^{-1}, \gamma_2^{-1} \alpha \gamma_2$. Since their determinant did not change, it is enough to show the result for one representative. We~take $p,p'$ as the base points of $\chi, \chi^{-1}$. Then we~know that $p$ is fixed as a set by $\alpha$ and $p'$ by $\beta$. Let $p_i,p_i'$ be the components of $p,p'$ over $\bar{\kk}$.
\begin{itemize}
\item We first assume that $p,p'$ are $3$\nobreakdash-points. We~now choose $\LL$-isomorphisms $\varphi, \varphi^{\op}$ from $S,S^{\op}$ to $\mathbb{P}^2$ as in Lemma \ref{3pLem}. Then we can see that $\varphi^{\op} \chi \varphi^{-1}$ has characteristic $2;1^3$,
\[
\Ind(\varphi^{\op} \chi \varphi^{-1}) = \{ [1\sep 0\sep 0],\, [0\sep 1\sep 0],\, [0\sep 0\sep 1]\},
\]
and the same for the inverse. Therefore, it is equal to $P D \sigma$ for some diagonal matrix $D$, $\sigma$ the standard Cremona involution and a permutation matrix $P$. Replacing $\varphi^{\op}$ with $P \circ \varphi^{\op}$ we may assume that $P = I_3$. Since $\Sigma:= (yz,xz,xy)$ satisfies the properties of an affine representative, we~conclude that $\sigma$ corresponds to a $\kk$-birational map from~$S$ to $S^{\op}$. We~therefore change our representative and choose $D=I_3$.

Since $p,p'$ are fixed as $3$\nobreakdash-points by $\alpha, \beta$, the affine representatives $A,B$ of $\alpha, \beta$ are diagonal matrices multiplied by permutation matrices. Those permutation matrices cannot be transpositions in our case. Otherwise, $\alpha$ would permute $p_1,p_2$ and fix $p_3$. Then for $g \in \Gal(\LL/\kk)$ that corresponds to a $3$\nobreakdash-cycle and sends $p_3$ to $p_1$, we~obtain
\[
g(p_3) = g(\alpha(p_3)) = \alpha(g(p_3)) = \alpha(p_1) = p_2,
\]
which is a contradiction.

In the case where the splitting field has degree $6$, it can also not correspond to a $3$\nobreakdash-cycle with the same type of argument. We~see in the case of a degree-$3$ splitting field $A_g, A_g^{\op}$ from Lemma \ref{3pLem} satisfy the condition to be representatives of $\kk$-automorphisms since for $A = A_g$ we have $A^g = A_g = A_g^{-1}A A_g$ and similarly for~$A_g^{\op}$. Therefore, we~can take $A = D_A A_g^{i_A}$, $B = D_B (A_g^{\op})^{i_B}$, where $0 \leq i_A,i_B \leq 2$ and $D_A,D_B$ are diagonal matrices. We~find there exists $\lambda \in \LL$ such that $\lambda B \Sigma = \Sigma A$, but since both $ B \Sigma $, $\Sigma A$ are affine representatives, we~have $\lambda \in \kk$, thus by changing $B$ we may assume $\lambda = 1$. Then
\[
D_B (A_g^{\op})^{i_B} \Sigma = B \Sigma = \Sigma A = \det(D_A) \xi^{i_A} D_A^{-1} (A_g^{\op})^{i_A} \Sigma
\]
(in affine coordinates). Then we~find
\[
\det(B) = \det(D_A)^2 \xi^{2i_A} = \det(A)^2 = \det(A)^{-1}
\]
in $\kk^*/(\kk^*)^3$.
\item If $p,p'$ are $6$\nobreakdash-points, we~can take $q$ as the fixed point of $\alpha$ (componentwise) from Lemma \ref{Fixpoints} and observe that $q' := \chi(q)$ is the componentwise fixed point of $\beta$.

We now choose $\LL$-isomorphisms $\varphi, \varphi^{\op}$ from $S,S^{\op}$ to $\mathbb{P}^2$ as in Lemma \ref{3pLem} for~$q,q'$. We~again write $A,B$ for the affine representatives of $\alpha, \beta$ and we write $Q$ for the affine representative of $\chi$. Since the image of $q$ under $\chi$ is fixed by $Q$, $Q$ fixes the points $[1\sep 0\sep 0]$, $[0\sep 1\sep 0]$, $[0\sep 0\sep 1]$, since $q' = \chi(q)$.

By Lemma $\ref{BaseP}$, $Q$ is of degree $5$. Since it fixes $[1\sep 0\sep 0]$, $[0\sep 1\sep 0]$, $[0\sep 0\sep 1]$ but cannot have them as base points, its first entry has an $x^5$ component but no $y^5,z^5$ components, and similarly for the second and third entries. Since $[1\sep 0\sep 0]$, $[0\sep 1\sep 0]$, $[0\sep 0\sep 1]$ are fixed by $A,B$ we conclude that $A,B$ are diagonal matrices. We~get another $\lambda \in \LL$ such that $\lambda B Q = Q A$ and we can eliminate it as in the first case. We~now obtain
\[
Q = B^{-1}Q A = B^{-1}Q(a_{11}x, a_{22}y, a_{33}z).
\]
Because of the $x^5$ term in the first entry we find $b_{11} = a_{11}^5$ and similarly $b_{22} = a_{22}^5, b_{33} = a_{33}^5$. Thus $B = A^5$ and we conclude that $\det(B) = \det(A)^5 = \det(A)^{-1}$ modulo $(\kk^*)^3$.\qedhere
\end{itemize}
\end{proof}
\begin{corollary}
For any Sarkisov link $\chi' \colon S \dra S^{\op}$ there exist $\alpha \in \Aut_{\kk}(S)$, $\beta \in \Aut_{\kk}(S^{\op})$ and $\chi \in R$ such that $\chi' = \beta \chi \alpha$ and the following expression $\det_R(\chi') := \det(\alpha) \det(\beta)^{-1}$ is well defined. For links $\chi' \colon S^{\op} \dra S$ we define $\det_R(\chi') := \det_R(\chi'^{-1})^{-1}$.
\end{corollary}

\begin{proof}
There is a $\chi \in R$, which has the same equivalence class as $\chi'$, which by definition gives the existence of $\alpha, \beta$. Assuming there exist $\alpha' \in \Aut_{\kk}(S)$, $\beta' \in \Aut_{\kk}(S^{\op})$ such that $\beta \chi \alpha = \chi' = \beta' \chi \alpha'$, then we conclude that $\chi = \beta'^{-1}\beta \chi \alpha \alpha'^{-1}$ and Lem\-ma~\ref{DetTrivRef} shows that $\det_R(\chi')$ is well defined.
\end{proof}

It remains to show that the elementary relations are sent onto the identity by our determinant map, \ie if $\tau_3\chi_3\tau_2\chi_2\tau_1\chi_1 = \id$ is an elementary relation between Sarkisov links of $S$, then we have $\prod_{i=1}^3 \det_R(\tau_i)\det_R(\chi_i) = 1$ in $\kk^* /(\kk^*)^3$.

\section{Rewriting elementary relations}
An elementary relation is a product of Sarkisov links $\chi_i \colon S \dra S^{\op}$, $\tau_i \colon S^{\op} \dra S$, $i=1,2,3$ based at $3$\nobreakdash-points such that $\chi_1,\chi_2,\chi_3$ are equivalent, $\tau_1,\tau_2,\tau_3$ are equivalent, and
\[
\tau_3\chi_3\tau_2\chi_2\tau_1\chi_1 = \id.
\]
By equivalence of Sarkisov links, there exist $\alpha_i, \delta_i \in \Aut_{\kk}(S)$, $\gamma_i, \beta_i \in \Aut_{\kk}(S^{\op})$ such that $\tau_i = \alpha_i \tau \beta_i$, $\chi_i = \gamma_i \chi \delta_i$, where $\chi, \tau$ are the links representing the respective equivalence class of links in $\mathcal{E}_3$. In~this section, we~prove that the choice of links representing the class does not affect the determinant of the word. In~Section \ref{Sec5} we then prove that the determinant is always equal to the class of $1$ in $\kk^*/(\kk^*)^3$.

\begin{lemma}\label{Subst}
Take an elementary relation
\[
\tau_3\chi_3\tau_2\chi_2\tau_1\chi_1 = \id
\]
between Sarkisov links of a Severi--Brauer surface $S$. Choose $\tau \colon S^{\op} \dra S$, $\chi \colon S \dra S^{\op}$ such that $\chi_i \sim \chi, \tau_i \sim \tau$. By equivalence of Sarkisov links, there exist $\alpha_i, \delta_i \in \Aut_{\kk}(S)$, $\gamma_i, \beta_i \in \Aut_{\kk}(S^{\op})$ such that
\[
\tau_i = \alpha_i \tau \beta_i, \quad \chi_i = \gamma_i \chi \delta_i.
\]
Then the product
\[
\prod_{i=1}^3 \det(\alpha_i\delta_i)\det(\beta_i\gamma_i)^{-1}
\]
does not depend on the choice of $\chi, \tau$.
In particular, the determinant of the word does not depend on the choice of representatives of the equivalence classes of links.
\end{lemma}

\begin{proof}
We take $ \tau' \colon S^{\op} \dra S$, $\chi' \colon S \dra S^{\op}$ such that $\tau \sim \tau'$, $\chi \sim \chi'$. Again, using the equivalence, we~find $\alpha, \delta \in \Aut_{\kk}(S)$, $\gamma, \beta \in \Aut_{\kk}(S^{\op})$ such that $\tau' = \alpha \tau \beta$, $\chi' = \gamma \chi \delta$ and find
\[
\tau_i = \alpha_i \alpha^{-1} \tau' \beta^{-1}\beta_i, \quad \chi_i = \gamma_i \gamma^{-1}\chi' \delta^{-1}\delta_i.
\]
The determinant of the word with $\tau, \chi$ is
\[
\prod_{i=1}^3 \det(\alpha_i\delta_i)\det(\beta_i\gamma_i)^{-1},
\]
and the one with $\tau', \chi'$ is
\begin{align*}
\prod_{i=1}^3 \det(\alpha_i \alpha^{-1}\delta_i\delta^{-1})\det(\beta_i\beta^{-1}\gamma_i\gamma^{-1})^{-1} &= \biggl(\prod_{i=1}^3 \det(\alpha_i\delta_i)\det(\beta_i\gamma_i)^{-1}\biggr) \Bigl(\frac{\det(\beta\gamma)}{\det(\alpha\delta)}\Bigr)^3
\\
&= \biggl(\prod_{i=1}^3 \det(\alpha_i\delta_i)\det(\beta_i\gamma_i)^{-1}\biggr)
\end{align*}
in $\kk^*/(\kk^*)^3$.
\end{proof}

\begin{lemma} \label{DepPoint}
Let
$\tau_3\chi_3\tau_2\chi_2\tau_1\chi_1 = \id$,
$\tau_3'\chi_3'\tau_2'\chi_2'\tau_1'\chi_1' = \id$,
be two elementary relations between Sarkisov links of a Severi--Brauer surface $S$. If $\chi_1 = \chi_1'$, $\Ind(\tau_1) = \Ind(\tau_1')$, then in $\kk^* /(\kk^*)^3$ we have the equality
\[
\prod_{i = 1}^3 \det_R(\chi_i) \det_R(\tau_i) = \prod_{i = 1}^3 \det_R(\chi_i') \det_R(\tau_i').
\]
\end{lemma}

\begin{proof}
Let $\chi := \chi_1 = \chi_1'$, $\tau:= \tau_1$. We~find $\alpha_i, \delta_i \in \Aut_{\kk}(S)$, $\gamma_i, \beta_i \in \Aut_{\kk}(S^{\op}) $ such that
\[
\tau_i = \alpha_i \tau \beta_i, \quad \chi_i = \gamma_i \chi \delta_i,
\]
and $\alpha_i', \delta_i' \in \Aut_{\kk}(S)$, $\gamma_i', \beta_i' \in \Aut_{\kk}(S^{\op})$ such that
\[
\tau_i' = \alpha_i' \tau \beta_i', \quad \chi_i' = \gamma_i' \chi \delta_i'.
\]
If $\chi \not \sim \tau^{-1}$ we may use $\tau^{-1}, \chi$ as the elements of $R$ that represent their respective equivalence of links. If $\chi \sim \tau^{-1}$ we can choose $\chi$ to represent its equivalence class and find that $\tau = \alpha \chi^{-1} \beta$ for some $\alpha \in \Aut_{\kk}(S)$ and $\beta \in \Aut_{\kk}(S^{\op})$. This implies that by definition $\det_R(\tau_i) = \det(\alpha_i \alpha) \det(\beta_i^{-1} \beta)$. In~the product of determinants from the statement, the determinants of $\alpha, \beta$ appear only in a power of $3$ and we may assume that they are the identity.

We take $p,p',q,q'$ as the base points of $\chi, \chi^{-1},\tau,\tau^{-1}$. Since $\Ind(\tau_1) = \Ind(\tau_1')$ and by our choice of representatives, we~have
\[
\alpha_1 = \delta_1 = \delta_1' =\id, \beta_1=\gamma_1=\beta_1'= \gamma_1' = \id.
\]
Due to the structure of the elementary relation, we~conclude that both $\delta_2, \delta_2'\alpha_1'$ send $\tau_1(p')$ to $p$. Therefore, there exists $\epsilon_1 \in \Aut_{\kk}(S)$ that fixes $p$ such that $\epsilon_1\delta_2 = \delta_2'\alpha_1'$. With Lemma \ref{DetTrivRef} we find $\epsilon_1' \!\in\! \Aut_{\kk}(S^{\op})$ such that $\epsilon_1' \chi \!=\! \chi \epsilon_1$ and $\det(\epsilon_1')\det(\epsilon_1) \!=\! 1$. We~therefore rewrite the second word as follows:
\begin{align*}
\tau_3'\chi_3'\tau_2'\chi_2'\tau_1'\chi &= \tau_3'\chi_3'\tau_2'(\gamma_2'\chi\delta_2')(\alpha_1'\tau_1) \chi = \tau_3'\chi_3'\tau_2'(\gamma_2'\chi)(\delta_2'\alpha_1')\tau_1 \chi
\\
&= \tau_3'\chi_3'\tau_2'(\gamma_2'\chi) (\epsilon_1\delta_2) \tau_1\chi = \tau_3'\chi_3'\tau_2'(\gamma_2'\chi \epsilon_1 \delta_2) \tau_1 \chi = \tau_3'\chi_3'\tau_2'(\gamma_2'\epsilon_1'\chi \delta_2)\tau_1\chi.
\end{align*}
We~can repeat this process and use this to find similar $\epsilon_i,\epsilon_i', i = 2, \dots ,5$ such that
\[
\beta'_2\gamma'_2 d'_1 = \epsilon_2 \beta_2 \gamma_2 , \quad \delta_3'\alpha_{2}'d'_{2} = \epsilon_{3} \delta_3\alpha_2, \quad \beta_3'\gamma_3' d'_{3} = \epsilon_{4} \beta_3 \gamma_3, \quad \alpha_3'd'_{4} = \epsilon_5 \alpha_3,
\]
and obtain
\[
\id = \tau_3'\chi_3'\tau_2'\chi_2'\tau_1'\chi = \epsilon_5\tau_3\chi_3\tau_2\chi_2\tau_1\chi = \epsilon_5.
\]
Hence, since $\det(\epsilon_i)\det(\epsilon_i')=1$ we can calculate that the determinants are the same.
\end{proof}
We note that
\[
\id = \tau_3\chi_3\tau_2\chi_2\tau_1\chi_1 = \alpha^{-1}\tau_3\chi_3\tau_2\chi_2\tau_1\chi_1 \alpha, \quad\alpha \in \Aut_{\kk}(S).
\]
Given the base points of $\chi_1^{-1}$ we have $\chi_1$ is unique up to right multiplication with an automorphism, therefore we can go even further and find that the determinant of the word $\tau_3 \chi_3 \tau_2\chi_2 \tau_1 \chi_1$ only depends, up to $(\kk^*)^3$, on the base points of $\chi_1^{-1}$ and $\tau_1$ which are both in $S^{\op}$.

\section{The determinant of an elementary relation} \label{Sec5}
If we fix the base points of $\tau_1$ and fix $\chi_1$ it is enough to choose one of such possible relations.
\[
\tau_3\chi_3\tau_2\chi_2\tau_1\chi_1 = \id
\]
and prove that $\prod_{i=1}^3 \det_R(\tau_i)\det_R(\chi_i) = 1$ in $\kk^* /(\kk^*)^3$. We~also choose $\chi:= \chi_1$ to represent its class of links. By Lemma \ref{DepPoint} this will imply the same type of result for all other relations with the same base points of $\tau_1$ and the same $\chi_1$.

In this chapter, we~will find, given $\tau$ and $\alpha_1 \in \Aut_{\kk}(S)$, suitable affine representatives $A_2, \dots , A_6$ of automorphisms $\alpha_2, \dots , \alpha_6$ such that
\begin{equation} \label{eq: exaRel}
\alpha_6 \tau \alpha_5 \chi \dots \tau \alpha_1 \chi = \id
\end{equation}
is an elementary relation. We~then take $\tau_i := \alpha_{2i} \tau \alpha_{2i-1}$, $\chi_i:= \chi$ and this will be the explicit elementary relation where we calculate the determinant of the word, \ie
\[
\prod_{i=1^3}\det(\alpha_{2i - 1}^{-1})\det(\alpha_{2i}).
\]

We start with the case where for the splitting field $\LL$ of the base point $p$ of $\chi$ we have $[\LL\sep \kk] = 3$. We~choose the conventional isomorphism $\varphi$ for $p$ and $\varphi^{\op}$ for the base point $p'$ of $\chi^{-1}$. We~have $\Sigma := (yz,xz,xy)$ is an affine representative of a $3$\nobreakdash-link equivalent to $\chi$. The affine representative of $\chi$ only differs from $\Sigma$ by a diagonal matrix, which is itself an affine representative of an automorphism, since $\Sigma$ is one. Thus, using Lemma \ref{Subst} we may assume that $\Sigma$ is an affine representative of $\chi$.
\begin{convention} \label{conv: 5.1}
For this section we take $\chi \colon S \dra S^{\op}$ a $3$\nobreakdash-link with base point $p$ whose splitting field $\LL$ satisfies $[\LL\sep \kk] = 3$ and has affine representative $\Sigma$. We~also take the $\varphi$ for $p$ and the $\varphi^{\op}$ for the base point $p'$ of $\chi^{-1}$. We~recall that we have $\xi$ such that
\[\arraycolsep4.5pt
A_g = \begin{pmatrix}
0 & 0 & \xi \\
1 & 0 & 0 \\
0 & 1 & 0
\end{pmatrix}\!, \qquad A_g^{\op} = \begin{pmatrix}
0 & 0 & \xi^{-1} \\
1 & 0 & 0 \\
0 & 1 & 0
\end{pmatrix}\!,
\]
where $[A_g] = \varphi \circ (\varphi^{-1})^g$, $[A_g^{\op}] = \varphi^{\op} \circ ((\varphi^{\op})^{-1})^g$. If there is no further specification we take all affine representatives with respect to $\varphi, \varphi^{\op}, A_g, A_g^{\op}$.
\end{convention}
Any elementary relation is constructed using two $3$\nobreakdash-points. We~will make a case distinction between the cases where they have the same class or where they do not have the same class.
\subsection{The relation with only one class of points}
We start by figuring out when exactly a relation of six links of a certain form is an elementary relation. Afterward, for this section, we will consider the case where the two $3$\nobreakdash-points in the elementary relation have the same class. Using the language of Lemma \ref{Subst} and supposing that $\tau, \chi$ are representatives of their respective equivalence classes of links, we~find $\tau^{-1} = \chi$.

\begin{lemma}\label{ElemRel}
Let $\Sigma_i$, $i=1,2,3$, be affine representatives of equivalent $3$\nobreakdash-links $\tau_i$ and let $\sigma_i$ be the projectivization of $\Sigma_i$. If
\begin{align*}
\sigma_i(\{[1\sep 0\sep 0],[0\sep 1\sep 0],[0\sep 0\sep 1] \}) &= \{[1\sep 0\sep 0],[0\sep 1\sep 0],[0\sep 0\sep 1] \},
\\
\sigma (\Ind(\sigma_{i}^{-1})) &= \Ind(\sigma_{i+1}), \quad i=1,2,
\\
\tag*{and}
\tau_3 \chi \tau_2 \chi \tau_1 \chi &= \id,
\end{align*}
then the last expression is an elementary relation.
\end{lemma}

\begin{proof}
We take $X_6$ the del Pezzo of degree $6$ that dominates $\chi$ and $X_{6,i}$ the del Pezzo of degree $6$ that dominates $\tau_i$. By Convention \ref{conv: 5.1} we see $p,p'$ are the base points of $\chi, \chi^{-1}$. Let $p_i,p_i'$ be the base points of $\tau_i, \tau_i^{-1}$. We~have $p_i \sim p_j$, $i,j = 1,2,3$ and $p'_i \sim p'_j$, $i,j = 1,2,3$. Since $X_{6,i}$ is the blow-up of $S$ in $p_i$, we~may assume that $X_{6,i} = X_{6,j}$, $i,j = 1,2,3$. We~denote this surface by $X_6'$. Take now $X_3$ the del Pezzo surface of degree $3$ that arises as the blow-up of $X_6$ in~$p_i$ or as the blow-up of $X_6'$ in~$p$. $X_3$ is a rank-$3$ fibration and therefore gives rise to elementary relations. We~want to prove that $\tau_3 \chi \tau_2 \chi \tau_1 \chi = \id$ is one of those relations.

Let $\rho_6 \dots \rho_1 = \id$ be such a relation where $\rho_i$ are Sarkisov links. Since the relation also depends on a blow-up model $X_3 \to S$, we~can choose the relation such that the base point of $\rho_1^{-1}$ is $p'$ and the base point of $\rho_2$ is $p_1$. We~therefore have $\rho_1 = \chi \alpha$ for some automorphism $\alpha$. Conjugating the relation with $\alpha$ we may assume that $\rho_1 = \chi$. Now in between links $\rho_{i + 1} \rho_i$ we can substitute an automorphism $\beta$ as $\rho_{i + 1} \beta \beta^{-1} \rho_i$ using the trivial relation. Therefore, we~can choose $\rho_2 = \tau_1$. We~now follow the diagram in Figure \href{https://jep.centre-mersenne.org/articles/10.5802/jep.347/}{1 in the original publication, p.1285} and prove that the Sarkisov links $\rho_i, i \geq 3$ are as in the statement.

By our assumptions, we~have $\chi(p_i') = p_{i + 1}$, $i = 1,2$ and $\tau_i(p') = p$. The link~$\rho_3$ is based at $\tau_2(p') = p$ therefore up to multiplication on the left side with an automorphism it is given by $\chi$. This automorphism on the left hand side can be substituted to $\rho_4$ using the trivial relation as stated before. Therefore, we~can choose $\rho_3 = \chi$. The base points of $\rho_4$ are $\chi(p_1') = p_2$, and we can substitute and choose $\rho_4 = \tau_2$. Since $\tau_2(p') = p$ is the base point of $\rho_5$, we~can, substituting again, assume that $\rho_5 = \chi$. The last link has the base point $\chi(p_2')$ and thus $\rho_6 = \beta \tau_6$ for some automorphism~$\beta$ of $S$. In~this case, we~cannot substitute anymore. We~see $\beta$ is chosen such that $\id = \beta \tau_6 \chi \dots \tau_1 \chi$. But the right hand side equals $\beta$ and thus $\rho_6 = \tau_3$.
\end{proof}
In the diagram of the elementary relation Figure \href{https://jep.centre-mersenne.org/articles/10.5802/jep.347/}{1 in the original publication, p.1285} we can see that the links act in a certain way on the base points of their neighboring links. To~create affine representatives of links that act the same, we~look at the action of $B \sigma A$ on the coordinate points for some matrices $A,B$ and $\sigma$ the standard Cremona involution.

\begin{lemma} \label{FixMatr}
Let $p,p'$ be as in Convention \ref{conv: 5.1}. Take $\rho \colon S^{\op} \dra S$ a $3$\nobreakdash-link with base points different from $p'$. Let $A = (a_1 \mid a_2 \mid a_3)$ be an affine representative of $\alpha \in \Aut_{\kk}(S^{\op})$ and $P$ an affine representative of $\rho$, whose existence follows from Proposition \ref{AffineRep}. Assume that $\alpha(p') \cup \Ind(\rho)$ are in general position.

Then $A^{\rho}D_{\xi}$ is an affine representative of an automorphism of $S$, where
\[
A^{\rho} = (P(a_1) \mid P(a_2) \mid P(a_3)), \quad D_{\xi} = \diag(1, \xi, \xi^2).
\]
Furthermore, its inverse $A' := D_{\xi}^{-1} (A^{\rho})^{-1}$ is an affine representative of $\alpha' \in \Aut_{\kk}(S)$, such that $\alpha' \rho \alpha$ sends $p'$ to $p$.
\end{lemma}

\begin{proof}
Let $a_i$, $i=1,2,3$ be the columns of $A$ and $b_i := P(a_i)$. Since
\[
A_g^{\op}(g(a_1) \mid g(a_2) \mid g(a_3)) = A_g^{\op} A^g = A A_g^{\op} = (a_2 \mid a_3 \mid \xi^{-1} a_1),
\]
we obtain
\begin{align*}
a_2 = A_g^{\op}(g(a_1)),\quad
a_3 &= A_g^{\op} (g(a_2)) = (A_g^{\op})^2(g^2(a_1)),\\
\xi A_g^{\op}(g(a_3)) &= \xi (A_g^{\op})^3 (a_1) = a_1.
\end{align*}
Since we have $A_g P^g (A_g^{\op})^{-1} = \xi P$, we~obtain
\begin{align*}
\xi b_2 = A_g(g(b_1)),
\quad
\xi b_3 &= A_g(g(b_2)) = \xi^{-1} A_g(g(\xi b_2)),
\\
\xi A_g(g(b_3)) &= \xi^{-1} A_g^3(b_1) = b_1.
\end{align*}
Therefore, if we look at
\[
A^{\rho} D_{\xi} = (b_1 \mid \xi b_2 \mid \xi^2 b_3) = (b_1 \mid A_g(g(b_1)) \mid A_g^2(g^2(b_1))),
\]
we observe that $A^{\rho} D_{\xi}$ satisfies the same properties as affine representatives and thus corresponds to an element in $\Aut_{\kk}(S)$.

We know that $[P] [A]$ sends $\{ [1\sep 0\sep 0], [0\sep 1\sep 0],[0\sep 0\sep 1]\}$ to the classes of $P(a_i)$ which are sent back to $\{ [1\sep 0\sep 0], [0\sep 1\sep 0],[0\sep 0\sep 1]\}$ by $A'$. Since we use $\varphi, \varphi^{\op}$ with $p,p'$, we~conclude that $\{ [1\sep 0\sep 0], [0\sep 1\sep 0],[0\sep 0\sep 1]\}$ corresponds to $p,p'$ and, therefore, $\alpha' \rho \alpha$ sends $p$ to $p'$, where $\alpha'$ is the automorphism represented by $A'$.
\end{proof}
\begin{remark}
In Lemma \ref{FixMatr} take any $\alpha'' \in \Aut_{\kk}(S^{\op})$ such that $\alpha'' \rho \alpha$ sends~$p$ to~$p'$. Then the affine representative of $\alpha''(\alpha')^{-1}$ is a diagonal matrix.
\end{remark}
Let $A_1$ be the affine representative of $\alpha_1$ in (\ref{eq: exaRel}). We~define $A_2 := (A_1^{\tau})^{-1}$, $A_3 := (A_2^{\chi})^{-1}, \dots , A_6 := (A_5^{\tau})^{-1}$, which are affine representatives of automorphisms $\alpha_i$ of $S$ if $i$ is even and of $S^{\op}$ if $i$ is odd (We remark that in this section $\tau = \chi^{-1}$). We~will prove that there exists an automorphism $\delta$ whose affine representative is a diagonal matrix such that replacing $\alpha_6$ by $\delta \circ \alpha_6$ yields an elementary relation as in (\ref{eq: exaRel}). We~will calculate that the determinant of this relation is $1$. By Lemma \ref{DepPoint} this will be enough to show that the determinant of any elementary relation where the classes of both $3$\nobreakdash-points are equal to the class of $p$ (\cf Convention \ref{conv: 5.1}) is equal to $1$.\enlargethispage{-4\baselineskip}

\begin{example}
We take an example from \cite[Cor.\,2.3.6]{BSY}, where $\kk := \mathbb{C}(u,v)$ and $\LL := \kk(\sqrt[3]{u})$. One can prove that this is a Galois extension of degree $3$ and that its norm map is non surjective, not having $\xi := v^{-1}$ in its image. Thus, using {\cite[Lem.\,2.3.3]{BSY}} there exists a Severi--Brauer surface $S_{\xi}$ and an isomorphism $\varphi \colon S_{\xi} \to \mathbb{P}_{\LL}^2 $ such that for $g \in \Gal(\LL/ \kk)$ a generator we have $\varphi \circ g \circ \varphi^{-1} = [A_g] \circ g$, where $A_g$ is as usual.

We now aim to find an example for the relation and check whether the determinant behaves like we want it to. For this, we~first try to find an $\alpha_1 \in \Aut_{\kk}(S_{\xi})$. For this, we~need an invertible matrix $A_1$ such that $A_gA_1^g = A_1A_g$. For simplicity we choose~$A_1$ to be over $\kk$, thus we need $A_gA_1 = A_1A_g$. To~do this, we~start with a column (or a first point) and use these equations to find the other column (its ``Galois orbit''). An~example of this would be
\[\arraycolsep4.5pt
A_1 := \begin{pmatrix}
1 & 3v & 2v \\
2 & 1 & 3v \\
3 & 2 & 1 \\
\end{pmatrix}\!.
\]
We now aim to find fitting $A_i$, $i=2, \dots , 6$ such that $A_6 \sigma A_5 \sigma \dots A_1 \sigma = \id$, where $\sigma$ is the standard Cremona involution. But we have more information on our relation. We~know that $A_i \sigma A_{i-1}$ always sends $\{ [1\sep 0\sep 0], [0\sep 1\sep 0], [0\sep 0\sep 1] \}$ to itself. We~assume that there is no permutation within the set and thus we can prove, for $a_{i}$ being the columns of $A_1$, that we can choose $(\Sigma(a_{1}) \mid \Sigma(a_{2}) \mid \Sigma(a_{3}) )^{-1}$ as $A_2$ (where we apply an affine version of $\sigma$ to the $a_i$). We~call this matrix $(A_1^{\Sigma})^{-1}$.

We can now find, since $\sigma \circ A_g = \xi A_g^{\op} \circ \sigma$, that if we instead choose $A_2 := D_{\xi^{-1}} (A_1^{\Sigma})^{-1}$, where $D_{\xi^{-1}}$ is the diagonal matrix with $1,\xi^{-1},\xi^{-2}$ on its diagonal, $A_2$ satisfies our restriction $A_g^{\op} A_2 = A_2 A_g^{\op}$. We~therefore continue with this principle and multiply with some fitting elements in $\kk$ we find the matrices:\vspace*{-3pt}
\bgroup
\begin{align*}
\arraycolsep4.5pt
A_2 &:= \begin{pmatrix}
36v^2 - 6 v & -3v & -18v + 4 \\
-18 v^2 + 4v & 36v^2 - 6 v & -3v \\
-3v^2 & -18 v^2 + 4v & 36v^2 - 6 v \\
\end{pmatrix}\!,\\
A_3 &:= \begin{pmatrix}
-9 v^2 + 2v & 6v^2 - v & 54v^3 - 21 v^2 + 2v \\
54v^2 - 21 v + 2 & -9 v^2 + 2v & 6v^2 - v \\
6v - 1 & 54v^2 - 21 v + 2 & -9 v^2 + 2v \\
\end{pmatrix}\!,\\
A_4 &:= \begin{pmatrix}
1 & 2 & 3 \\
3v & 1 & 2 \\
2t-2 & 3v & 1 \\
\end{pmatrix}\!,\\
A_5 &:= \begin{pmatrix}
36v^2 - 6 v & -18 v^2 + 4v & -3v^2 \\
-3t-2 & 36v^2 - 6 v & -18 v^2 + 4v \\
-18 v + 4 & -3v & 36v^2 - 6 v \\
\end{pmatrix}\!,\\
A_6 &:= \begin{pmatrix}
-9 v^2 + 2v & 54v^2 - 21 v + 2 & 6v - 1 \\
6v^2 - v & -9 v^2 + 2v & 54v^2 - 21 v + 2 \\
54v^3 - 21 v^2 + 2v & 6v^2 - v & -9 v^2 + 2v \\
\end{pmatrix}\!,
\end{align*}
\egroup%
and get $\det(A_1A_3A_5)\det(A_2A_4A_6)^{-1} = 1$. We~still do not know what\vspace*{-3pt}
\[
A_6 \sigma A_5 \sigma \dots A_1 \sigma
\]
is, we~just know that it is a diagonal matrix. But one can calculate this with a computation software and find that it is the identity. Therefore, we~have seen an example of an elementary relation, where the two $3$\nobreakdash-points that give the elementary relation have the same equivalence class. Moreover, the determinant of the relation is~$1$.
\end{example}

\begin{remark} \label{rem: image under sigma}
In Lemma \ref{LargeCalc} we need to study the image of certain lines and conics under $\sigma$. To~this end, we~will now study how to calculate them.

The map $\sigma = \pi_2 \circ \pi_1^{-1}$, where $\pi_1$ is the blow up of the three points $[1\sep 0\sep 0]$, $[0\sep 1\sep 0]$, $[0\sep 0\sep 1]$ and $\pi_2$ the contraction of the strict transforms of the lines through respectively two out of the three points.

Take $L$ a line such that $L$ contains none of the three base points of $\sigma$. The strict transform of $L$ with respect to $\pi_1$ has self intersection $1$ since it does not contain the base points. The line $L$ intersects the three lines contracted in $\pi_2$ each exactly once, and therefore its image has a self intersection $4$ and contains the three base points of~$\sigma$. This makes it a conic through $[1\sep 0\sep 0]$, $[0\sep 1\sep 0]$, $[0\sep 0\sep 1]$. Performing the exact inverse calculation one can prove that the image of a conic through those three points is a line not containing any base point of $\sigma$.

Take $C$ a conic through exactly two of the base points $p_i,p_j$ of $\sigma$. The strict transform of $C$ under $\pi_1$ has self intersection $2$ and does not intersect the line through $p_i,p_j$. It intersects the other two lines contracted by $\pi_2$ exactly once and therefore the image of $C$ has self-intersection $4$ and contains $p_i,p_j$. Therefore, it is a conic through $p_i,p_j$.
\end{remark}

\begin{lemma} \label{LargeCalc}
Let $\mathbf{U}$ be a field and $A \in \GL_3(\mathbf{U})$ ( where the classes of its columns in $\mathbb{P}^2_{\emph{C}}$ together with the coordinate points are in general position). Then if we define $M_1(A):= A$, $M_i(A) := (M_{i-1}(A)^{\Sigma})^{-1}$, $i=2, \dots , 6$, where $\Sigma = (yz,xz,xy)$, we~get
\begin{align*}
d(A)&:= \det(A^{-1}M_3(A)^{-1}M_5(A)^{-1})\det(M_2(A)M_4(A)M_6(A))
\\
&\hphantom{:}= \frac{P(A)^{14}}{P(A^{-1})^7 \det((A^{\Sigma}))^{42}\det(A)^{42}},
\end{align*}
where $P$ is the product of all entries in the matrix. And as birational maps over $\mathbb{P}^2$ we get $M_6(A) \sigma \dots M_2(A) \sigma A \sigma = D(A)$, where $D(A)$ is a diagonal matrix such that:
\[
D(A)_{ii} = \det(A)^6\frac{\prod_{k=1}^3 b_{ik}^2}{\prod_{l=1}^3a_{li}},
\]
where $A = (a_{ij})_{i,j = 1}^3, A^{-1} = (b_{ij})_{i,j = 1}^3$. Thus $\det(D(A)) =\sfrac{P(A^{-1})^2\det(A)^{18}}{P(A)}$.
\end{lemma}

\begin{proof}
To calculate the determinant, we~need to calculate $\det(M_i(A))$, $i=2, \dots , 6$. We~calculate directly:
\[
\det(M_3(A))=-(P(A)\det(A)\det(M_2(A))^4)^{-1}.
\]
With this we can calculate $\det(M_i(A)), i \geq 3$ using $M_i(A) = M_{i-1}(M_2(A))$ and replacing $\det(M_{i-1}(\cdot))$ with the previous equality
\begin{align*}
\det(M_4(A))&=-P(M_2(A))^{-1}P(A)^4\det(A)^4\det(M_2(A))^{15},
\\
\det(M_5(A))&=P(M_3(A))^{-1}P(M_2(A))^4P(A)^{-15}\det(A)^{-15}\det(M_2(A))^{-56},
\\
\det(M_6(A))&=P(M_4(A))^{-1}P(M_3(A))^4P(M_2(A))^{-15}P(A)^{56}\det(A)^{56}\det(M_2(A))^{209}.
\end{align*}
Thus, we~need to find $P(M_i(A))$, $i=2,3,4$:
\[
P(M_2(A)) = -P(A^{-1})P(A)^2\det(A)^9\det(M_2(A))^9.
\]
Using $P(M_2(A)^{-1}) = P(A)^2$ and $M_i(A) = M_{i-1}(M_2(A))$ we obtain
\begin{align*}
P(M_3(A)) &= P(A)^{-3}P(A^{-1})^2\det(A)^9\det(M_2(A))^{-9},
\\
P(M_4(A))&=P(A^{-1})^{-3}P(A)^7\det(A)^{-18}\det(M_2(A))^{18}.
\end{align*}
Combining these two equations, we~obtain
\begin{multline*}
\det(A^{-1}M_3(A)^{-1}M_5(A)^{-1})\det(M_2(A)M_4(A)M_6(A))
\\
= \frac{P(A)^{14}}{P(A^{-1})^7 \det((A^{\Sigma}))^{42}\det(A)^{42}}.
\end{multline*}
The fact that $M_6(A) \sigma \dots M_1(A) \sigma$ is a diagonal automorphism can be established using Remark \ref{rem: HomSys} to study the characteristics and base points of
\[
M_6(A) \sigma \dots \sigma M_i(A),\quad i = 1, \dots , 5.
\]
Take $P := \{[1\sep 0\sep 0],[0\sep 1\sep 0],[0\sep 0\sep 1]\}$
\begin{enumerate}
\item For $i = 5$ we have the map has characteristic $2;1^3$ with the three base points $M_5(A)^{-1}(P)$. In~addition, it sends each point in $P$ to itself.
\item For $i = 4$ we have the map has characteristic $4;2^3,1^3$ with the three base points $M_4(A)^{-1}(P)$ of multiplicity $2$ and the three base points
\[
M_4(A)^{-1}(\sigma(M_5(A)^{-1}(P))) = P
\]
of multiplicity $1$.
\item For $i = 3$ we have the map has characteristic $5;2^6$ with the six base points
\[
M_3(A)^{-1}(\sigma(M_4(A)^{-1}(P))) = P, \quad M_3^{-1}(P).
\]
\item For $i = 2$ we have the map has characteristic $4;2^3,1^3$ with the three base points $M_2(A)^{-1}(P)$ of multiplicity $1$ and the three base points
\[
M_2(A)^{-1}(\sigma(M_3(A)^{-1}(P))) = P
\]
of multiplicity $2$.
\item For $i = 1$ we have the map has the characteristic $2;1^3$ with the base points
\[
M_1(A)^{-1}(\sigma(M_2(A)^{-1}(P))) = P.
\]
\end{enumerate}
By the same argument applied to the inverse, the inverse of the map at $i = 1$ has characteristic $2;1^3$ with base points $P$. Therefore, the map $M_6(A) \sigma \dots M_1(A) \sigma$ has degree $1$, \ie is an automorphism and maps $P$ to $P$. It remains to prove that the points are not permuted. It suffices to prove that the map at $i = 1$ contracts the line~$L$ through $p_i, p_j \in P, i \neq j$ to the point in $P$ that is not contained in the line. To~this end, we~study the image of $L$ by the map $M_6(A) \sigma \dots \sigma M_1(A)$. To~calculate the image of a curve under $\sigma$, we~refer to Remark \ref{rem: image under sigma}.
\begin{enumerate}
\item We find $M_2(A) \sigma M_1(A)$ fixes the components of $P$ and has base points
\[
M_1(A)^{-1}(P)
\]
that by the general position assumption do not lie on the line. The image of $L$ is the conic through $M_2(A)(P)$ and $p_i,p_j$.
\item The image of this conic under $\sigma$ is the conic through $\sigma(M_2(A)(P)), p_i, p_j$.
\item We note that $M_3(A)(\sigma(M_2(A)(P))) = P$ and $M_4(A) \sigma M_3(A)$ fix $P$. Therefore, the image under $M_4(A) \sigma M_3(A)$ of the conic through
\[
\sigma(M_2(A)(P)), p_i, p_j
\]
is the line through $p_i,p_j$.
\item The map $\sigma$ contracts this line onto the point not contained in $L$ which is then fixed by $M_6(A) \sigma M_5(A)$.
\end{enumerate}

A similar argument as above just using the inverse relation that this diagonal must map the classes of the columns of $(A^{-1})^{\Sigma}$ to the classes of the columns of $M_6(A)$. Thus, we~now calculate the matrices $M_i(A)$ and confirm the form of the matrix $D(A)$ as given in the statement.

We have for $\lambda_i(A) \in \kk^*$, $i = 2, \dots , 6$:
\[
M_2(A)_{ij} = (\prod_{k \neq i} a_{jk}) b_{ij} \lambda_2(A),
\quad
\lambda_2(A) = \frac{\det(A)}{\det(A)^{\Sigma}}.
\]
Using $M_i(A) = M_{i-1}(M_2(A))$ and dividing by $P(A), P(A^{-1})$ if needed, we~obtain
\begin{align*}
M_3(A)_{ij} &= \frac{\prod_{k \neq i}b_{jk}}{\prod_{k = 1}^3a_{ik}} \lambda_3(A),&
\lambda_3(A) &= \lambda_2(M_2(A)) \lambda_2(A)^2P(A),
\\
M_4(A)_{ij} &= \frac{1}{(\prod_{k \neq i}a_{jk})(\prod_{k=1}^3 b_{ik})} \lambda_4(A),&
\lambda_4(A) &= \frac{\lambda_3(M_2(A))}{\lambda_2(A)},
\\
M_5(A)_{ij} &= \frac{\prod_{k=1}^3a_{ik}}{\prod_{k \neq i}b_{jk} \prod_{k \neq i} \prod_{l \neq j} a_{kl}} \lambda_5(A),&
\lambda_5(A) &= \frac{\lambda_4(M_2(A))}{P(A) \lambda_2(A)^2},
\\
M_6(A)_{ij} &= \frac{(\prod_{k=1}^3b_{ik}) (\prod_{k \neq i}a_{jk}^2)a_{ji}}{(\prod_{l \neq j}a_{li}) (\prod_{k \neq i} \prod_{l \neq j}b_{kl})} \lambda_6(A),&
\lambda_6(A) &= \frac{\lambda_5(M_2(A)) \lambda_2(A)}{P(A)}.
\end{align*}
To prove the form of $D(A)$ we define:
\[
D_1(A)_{ii} := \det(A)^6\frac{\prod_{k=1}^3b_{ik}^2}{\prod_{l=1}^3a_{li}},
\quad
(A^{-1})^{\Sigma}_{ij} = \prod_{l \neq i} b_{lj}.
\]
Thus:
\begin{align*}
(D_1(A)(A^{-1})^{\Sigma})_{ij} &= \det(A)^6\frac{\prod_{k=1}^3b_{ik}^2\prod_{l \neq i} b_{lj}}{\prod_{l=1}^3a_{li}}
\\
\tag*{and}
M_6(A)_{ij} &= \frac{(\prod_{k=1}^3(A^{-1})^2_{ik}) (\prod_{k = 1}^3a_{jk}^2)(\prod_{l \neq i}b_{lj})}{(\prod_{l = 1}^3a_{li})} \frac{\lambda_6(A)}{P(A^{-1})}.
\end{align*}
Therefore, for the diagonal matrix $D_2(A)$ with
\[
D_2(A)_{jj} = \biggl(\prod_{k=1}^3a_{jk}^2\biggr) \frac{\lambda_6(A)}{P(A^{-1})\det(A)^6},
\]
we obtain
\[
M_6(A) = D_1(A) (A^{-1})^{\Sigma}D_2(A).
\]
Thus, the diagonal $D_1(A)$ sends the projective class of the columns of $(A^{-1})^{\Sigma}$ to the ones of the columns of $M_6(A)$. This is the same action as $D(A)$ and since they are diagonal and the columns of $M_6(A)$ are in general position to the coordinate points, we~have $D(A) = D_1(A)$.
\end{proof}

\begin{remark} \label{RemDiag}
The matrix $M_2(A)$ is up to a diagonal, the unique matrix $M$ such that $M \sigma A$ fixes the coordinate points.

Let $D_i$, $i=2, \dots ,6$ be diagonal matrices and define
\[
M'_1(A) := A, \quad M'_i(A) := D_i M_2(M'_{i-1}(A)), \quad i = 2, \dots , 6.
\]
We study what happens to the product and the determinant if we use $M_i'$ instead of~$M_i$. For any diagonal $D$
\[
M_2(DM_i(A)) = \det(D)^{-1}M_2(M_i(A))D, \quad M_2(M_i(A)D) = D^{-2}M_2(M_i(A)).
\]
And therefore we have in projective coordinates
\begin{align*}
M'_3(A) &= \lambda_3 D_3 M_3(A) D_2, & \lambda_3 &= \det(D_2)^{-1},
\\
M'_4(A) &= \lambda_4 D_4 D_2^{-2} M_4(A)D_3, &\lambda_4 &= \det (D_2)^2 \det(D_3)^{-1},
\\
M'_5(A) &= \lambda_5 D_5D_3^{-2}M_5(A)D_4D_2^{-2}, &\lambda_5 &= \det(D_2)^{-2} \det(D_3)^2 \det(D_4)^{-1},
\\
M'_6(A) &= \lambda_6 D_6 D_4^{-2}D_2^4 M_6(A)D_5D_3^{-2}, &\lambda_6&= \det(D_2)^{4} \det(D_3)^{-2} \det(D_4)^{2} \det(D_5).
\end{align*}
Using $[D] \sigma [D] = \sigma$ we have, in projective coordinates:
\[
M'_6(A) \sigma \dots \sigma M'_1(A)\sigma = D_6 D_4^{-2}D_2^4 M_6(A) \sigma \dotsm \sigma A \sigma = D_6 D_4^{-2}D_2^4 D(A)
\]
as birational maps. One can also calculate:
\begin{align*}
\det(M'_1(A)^{-1}&M'_3(A)^{-1}M'_5(A)^{-1})\det(M'_2(A)M'_4(A)M'_6(A))
\\
&= d(A) \det(D_2)^{31}\det(D_3)^{-15} \det(D_4)^{7}\det(D_5)^{-3} \det(D_6)
\\
&= \det(D_6 D_4^{-2}D_2^4) d(A) (\det(D_2)^{9}\det(D_3)^{-5} \det(D_4)^{3}\det(D_5)^{-1})^3.
\end{align*}
\end{remark}
Using the preceding results, we~now give our first example of an elementary relation.

\begin{lemma} \label{TrivEx}
Let $A_1$ be an affine representative of a $\kk$-automorphism of $S^{\op}$, where the classes of the columns of $A_1$ in $\mathbb{P}^2_{\LL}$ together with the coordinate points are in general position. We~define with the notation of Lemma \ref{LargeCalc}:
\[
A_i := D_{\xi^{t(i)}}M_2(A_{i-1}), \; i=2, \dots , 5 , \quad A_6 := D(A_1)^{-1}D_{\xi}^3D_{\xi^{t(6)}}M_2(A_5),
\]
where $t(i) = -1$, $i \text{ even }$, $t(i) = 1 $, $i \text{ odd}$. Then $A_i$ are affine representatives of automorphisms $\alpha_i$ and $\alpha_6 \chi \dots \alpha_1 \chi = \id$ and this is an elementary relation whose determinant is in $(\kk^*)^3$.
\end{lemma}

\begin{proof}
We use Remark \ref{RemDiag} with
\[
D_i = D_{\xi^{t(i)}},\; i = 2, \dots , 5\quand D_6 = D(A_1)^{-1}D_{\xi}^3D_{\xi^{t(6)}}.
\]
Thus in projective coordinates we have $A_6 \sigma \dots A_1 \sigma = D_6 D_4^{-2}D_2^4 D(A_1) = \id$.

The fact that the $A_i$ are affine representatives follows from Lemma \ref{FixMatr} which implies that $D_{\xi^{t(i)}}M_2(A_{i-1})$ is an affine representative and we can calculate that $D(A_1)^{-1}D_{\xi}^3$ is one as well. In~Lemma \ref{LargeCalc}, we~can see that $d(A_1)$ equals $\det(D(A_1))$ up to $(\kk^*)^3$. Therefore, Lemma \ref{LargeCalc} and Remark \ref{RemDiag} imply that
\[
\det(A_2A_4A_6)\det(A_1^{-1}A_3^{-1}A_5^{-1}) = \det(D_{\xi})^{-54} d(A_1) \det(D(A_1))^{-1} \in (\kk^*)^3,
\]
since by Lemma \ref{matrixL} we see $P(A_1) \in \kk^*$. The fact that this relation is elementary follows from the action of $A_i$ on the base points and Lemma \ref{ElemRel}.
\end{proof}
Using this elementary relation, we~can now prove the first case of Theorem \ref{MainA} where $\tau^{-1} = \chi$ and where $\LL_p$ has degree $3$ over $\kk$.

\begin{theorem} \label{TrivialCase}
All elementary relations between Sarkisov links of $S$ where both classes of base points are equivalent and their splitting field is a degree-$3$ extension of~$\kk$ have a trivial determinant.
\end{theorem}

\begin{proof}
We fix an affine representative $A_1$ (where the classes of the columns of $A_1$ in $\mathbb{P}^2_{\LL}$ are in general positions) of a $\kk$-automorphism of $S$. By Corollary \ref{DepPoint}, it suffices to prove the theorem for one possible choice for the other $5$ automorphism. In~this case, we~also have only one kind of link for which we will use the one represented by~$\Sigma$. We~choose affine representatives $A_i$ as given in Lemma \ref{TrivEx}, which implies that the determinant is trivial.
\end{proof}

\subsection{The relation with two classes of $3$\nobreakdash-points}
If the base points of $\tau,\chi^{-1}$ are not equivalent, we~again reduce this problem to the case of Lemma \ref{LargeCalc}. The difficulty is that $\tau^{-1} \neq \chi$ and thus we need some information on the base points of $\tau$ and its inverse. The next few results identify these points, an affine representative of $\tau$ and the $\kk$-automorphisms to find a relation.

\begin{lemma}
Let $p,q \in S$ be two non-equivalent $3$\nobreakdash-points with splitting fields $\LL_p, \LL_q$ such that $ [\LL_p\sep \kk] = [\LL_q\sep \kk]=3$. Let $\FF$ be the compositum of $\LL_p, \LL_q$. Then we find $\Gal(\FF/\kk) \simeq \Gal(\FF/\LL_p) \times \Gal(\FF/\LL_q)$
\end{lemma}

\begin{proof}
We recall that $\kk \subset \LL_P, \LL_q$ are Galois extensions. We~take $\omega_p, \omega_q \in \Bar{\kk}$ such that $\LL_p = \kk(\omega_p), \LL_q = \kk(\omega_q)$ using the Primitive Element Theorem. Then we find $\FF = \kk(\omega_p,\omega_q)$. Since $p,q$ are not equivalent, we~conclude that $\LL_p \cap \LL_q = \kk$. Assume that a zero of the minimal polynomial of $\omega_p$ over $\kk$ lies in $\LL_q$, then by normality of $\LL_q$ we have $\LL_p \subset \LL_q$, and hence $\LL_p = \LL_q$, a contradiction. Thus we have $\LL_q \subset \FF = \LL_q(\omega_p)$ is of degree $3$ and therefore $\kk \subset \FF$ is of degree $9$. Then the map
\begin{align*}
\Gal(\FF/\kk) &\to \Gal(\FF/\LL_p) \times \Gal(\FF/\LL_q)
\\
g &\mto (g_{|\LL_p}, g_{|\LL_q})
\end{align*}
is injective, since $\omega_p, \omega_q$ generate $\FF$ and thus since both sides contain $9$ elements, it is also surjective, and thus it is an isomorphism.
\end{proof}

\begin{convention} \label{con:2SF}
From now on we write $\LL_p, \LL_q$ for the splitting fields of non-equivalent $3$\nobreakdash-points $p,q$ of $S$ and $\FF := \LL_p\LL_q$. We~always will identify $\Gal(\FF/\kk) $ as $\Gal(\FF/\LL_p) \times \Gal(\FF/\LL_q)$. We~will also extend our isomorphisms $\varphi$ as in Lemma \ref{3pLem} (for $p$) to $\FF$. Since this can still be represented over $\LL_p$, the $\LL_q$ part of the Galois group of $\FF$ does not affect $\varphi$. In~the case where $\LL_q$ is a degree-$3$ extension we take $g'$ a generator of $\Gal(\LL_q/\kk)$.
\end{convention}
We extend the preceding constructions to our larger field $\FF$. To~this end, we~extend our isomorphism $\varphi$ to $\FF$ and then study the new twisted action on this. With the new twisted action, we~can retain our notion of affine representatives while working with~$\FF$.

\begin{lemma}
Let $\kk$, $\LL_p$, $\LL_q$, $\FF, \varphi$ be as in Convention \ref{con:2SF} where $\kk \subset \LL_p$ is of degree $3$. For $A_g$ as usual, we~conclude that for every $(g,u) \in \Gal(\FF/\kk) \simeq \Gal(\FF/\LL_p) \times \Gal(\FF/\LL_q)$
\[
\varphi \circ (g^i,u) \circ \varphi^{-1} = [(A_g)^i] \circ (g^i,u), \quad i = 0,1,2.
\]
Furthermore, for each $3$\nobreakdash-point $q'$ equivalent to $q$, the components of $\varphi(q')$ are fixed by $[A_g] \circ (g,\id)$ and permuted by $\{ \id \} \times \Aut_{\kk}(\LL_q)$, in the same way as $q'$ in $S_{\FF}$.

The properties of Theorem \ref{AffineRep} also remain valid for the extension $\kk \subset \FF$, where we replace $g^{i}$ with $(g^{i},u)$ and choose $A_{(g^i,u)} = A_{g^i}$.
\end{lemma}

\begin{proof}
We take $\varphi$ as we did before for $\LL_p$, since it can be defined over $\FF$, and we get
\[
\varphi \circ (g^i,u) \circ \varphi^{-1} \circ (g^i,u)^{-1} = \varphi \circ (\varphi^{-1})^{(g^i,u)} = \varphi \circ (\varphi^{-1})^{g^i} = [(A_g)^i].
\]
Since $(g,\id)$ fixes the geometric points $q_i'$ on $S_{\FF}$ we have $[A_g] \circ (g,\id)$ fixes $\varphi(q'_i)$. Since $(\id,u)$ permutes the $q_i'$ and $\varphi \circ (\id,u) \circ \varphi^{-1} = (\id,u)$, the same holds for $(\id,u)$ with $\varphi(q_i')$.

Since the elements in $\varphi \Bir_{\kk}(S) \varphi^{-1}$ can be defined over $\LL_p$, we~may choose the same affine representation as in Theorem \ref{AffineRep}.
\end{proof}

We now want to study which points in $\mathbb{P}^2_{\FF}$ correspond to $3$\nobreakdash-points $q'$ that split over~$\LL_q$. By the preceding lemma, these points must be fixed by $A_g \circ (g,\id)$. Assume that such a $q'$ has a component on one of the geometric lines between the components of $p$ in $S$. Then all the components of $q'$ would lie on that line, since the lines are fixed by $\Aut_{\kk}(\LL_q)$. But then $g$ would permute the lines and fix $q'$ which is impossible.

If a geometric conic existed through all six components of those two $3$\nobreakdash-points its strict transform under the link $\chi$ based at $p$ would be a geometric line through $\chi(q')$ which is still a $3$\nobreakdash-point and hence this is impossible. Therefore, $p \cup q'$ are in general position.
\begin{remark}
In Lemma \ref{3-pointNonEq} we use the field norm $N_{\FF/\LL_q}$. The Galois group $\Gal(\FF/\LL_q)$ is generated by $(g,\id)$ (which we abbreviate to $g$) and therefore we have $N_{\FF/\LL_q}(a) = a g(a)g^2(a)$, $a \in \FF$.
\end{remark}

\begin{lemma} \label{3-pointNonEq}
Let $p,q,\LL_p,\LL_q,g'$ be as in Convention \ref{con:2SF} and assume that $\kk \subset \LL_p, \LL_q$ are of degree $3$. We~can construct the following representations of $q$ over $\FF = \LL_p \LL_q$. There exists $b \in \FF$ such that $N_{\FF / \LL_q}(b) = \xi^{-2}$ and for each such $b$, define
\[\arraycolsep4.5pt
B_S(b) := \begin{pmatrix}
1 & 1 & 1 \\
\sfrac{1}{\xi g(b)} & \sfrac{1}{\xi g(g'(b))} & \sfrac{1}{\xi g(g'(g'(b)))} \\
b & g'(b) & g'(g'(b))
\end{pmatrix}\!.
\]
The classes of the columns of $B_S(b)$ form a $3$\nobreakdash-point over $S$ via the usual isomorphism $\varphi$ (extended as in Convention \ref{con:2SF}) and each $3$\nobreakdash-point equivalent to $q$ not on the coordinate lines can be given like this. We~also have
\[
A_gB_S(b)^g = \xi B_S(b) D_b^g, \quad D_b = \diag(b, g'(b),g'(g'(b))).
\]
Furthermore,
\[\arraycolsep4.5pt
B_S(b)^{g'} = B_S(b) \begin{pmatrix}
0 & 0 & 1 \\
1 & 0 & 0 \\
0 & 1 & 0
\end{pmatrix}\!.
\]
\end{lemma}

We will later refer to the permutation matrix used in the above formula as $(123)$).

\begin{proof}
Consider $[1\sep b_1\sep b_2]$ satisfying
\[
[1\sep \xi^{-1}g(b_2^{-1})\sep  \xi^{-1}g(\sfrac{b_1}{b_2})] = [\xi g(b_2)\sep 1\sep g(b_1)] \stackrel{!}{=} [A_g](g([1\sep b_1\sep b_2])) = [1\sep b_1\sep b_2],
\]
from which it follows that $b_1 = g(b_2^{-1})\xi^{-1}$ and $N_{\FF/\LL_q}(b_2) = \xi^{-2} $. Since such a $3$\nobreakdash-point is a $g'$-orbit, this gives the stated structure of $B_S(b)$.

Since $q$ exists and does not lie on one of the lines between the components of $p$, such a $b$ exists and therefore also of $B_S(b)$.

Conversely, for every $b \in \FF$ such that $N_{\FF / \LL_q}(b) = \xi^{-2}$ we get
\[
[A_g](g([1\sep \sfrac{1}{\xi g(b)}\sep b])) = [\xi g(b)\sep 1\sep \sfrac{1}{\xi g(g(b))}] = [1\sep \sfrac{1}{\xi g(b)}\sep  b].
\]
Thus, the columns of $B_S(b)$ are fixed by $A_g \circ g$. Since the columns of $B_S(b)$ form an orbit with respect to $\Gal(\LL_q/\kk)$ we get, with the previous property, that the columns form a $3$\nobreakdash-point.

The identity $A_gB_S(b)^g = \xi B_S(b) D_b^g$ follows by direct calculation using
\[
[A_g](g([1\sep \sfrac{1}{\xi g(b)}\sep b])) = [\xi g(b)\sep 1\sep \sfrac{1}{\xi g(g(b))}] = [1\sep \sfrac{1}{\xi g(b)}\sep  b].
\]
By the construction of $B_S(b)$ we can directly find
\[\arraycolsep4.5pt
B_S(b)^{g'} = B_S(b) \begin{pmatrix}
0 & 0 & 1 \\
1 & 0 & 0 \\
0 & 1 & 0
\end{pmatrix}\!.\qedhere
\]
\end{proof}

Given such matrices, we~find an affine representative of a Sarkisov link $\tau$ that has $q$ as base points if the splitting field of $q$ is a degree-$3$ extension. Since the determinant of the word depends only on the base points of $\chi_1^{-1}$ and $\tau$ this type of link suffices as a choice of $\tau_1$.

\begin{lemma} \label{BB'}
Take $\LL_p, \LL_q, \FF, \varphi$ as in Convention \ref{con:2SF}, where $\kk \subset \LL_p, \LL_q$ are of degree $3$. We~can take $b \in \FF$, so that $N_{\FF / \LL_q}(b) = \xi^{2}$ and with Lemma \ref{3-pointNonEq} we define $B := B_{S^{\op}}(b)$ (representing a $3$\nobreakdash-point of $S^{\op}$) and $B':=B_S(b^{-1})$ (representing a $3$\nobreakdash-point of $S$).
Then, there is a $\tilde \lambda \in L_p$ such that
\[
\tilde \Sigma := (\tilde \lambda B') \circ \Sigma \circ B^{-1}
\]
is an affine representative of a $3$\nobreakdash-link from $S^{\op}$ to $S$, where $\Sigma = (yz,xz,xy)$.
\end{lemma}

We~will later refer to the projectivization of $\tilde \Sigma$ as $\tilde \sigma$.

\begin{proof}
The existence of $b$ follows from Lemma \ref{3-pointNonEq}, where in this case we need $N_{\FF/\LL_q}(b) = \xi^2$ instead of $\xi^{-2}$ since we work over $S^{\op}$ instead of $S$. Since
\[
N_{\LL_p / \kk}(\det(D_b)) = N_{\LL_p / \kk}(N_{\FF/\LL_p}(b)) = \xi^6,
\]
we have $N_{\LL_p / \kk}(\sfrac{\xi^2}{g(\det(D_b))}) = 1$ and therefore, using Hilbert's Theorem 90, we~find a $\tilde \lambda \in L_p$ such that
\[
\frac{\tilde \lambda}{g(\tilde \lambda)} = \frac{\xi^2}{g(\det(D_b))}.
\]
We now get\vspace*{-3pt}
\begin{align*}
(\tilde \Sigma)^{g'} &= (\tilde \lambda B'^{g'}) \circ \Sigma \circ (B^{-1})^{g'}
\\
&= (\tilde \lambda B' (123)) \circ \Sigma \circ ((123)^{-1} B^{-1}) = (\tilde \lambda B' (123) (123)^{-1}) \circ \Sigma \circ B^{-1} = \tilde \Sigma.
\end{align*}
This shows that $\tilde \sigma$ is defined over $\LL_p$. To~prove the properties of affine representatives (since $\tilde \sigma$ is surely $\FF$-birational, it is also sufficient) we look at:\vspace*{-3pt}
\begin{align*}
\xi^{-1}(\tilde \Sigma)^g &= (g(\tilde \lambda)\xi^{-1} B'^g) \circ \Sigma \circ (B^{-1})^g = (g(\tilde \lambda) A_g^{-1} B' D_{b^{-1}}^g) \circ \Sigma \circ (\xi (D_{b}^{-1})^g B^{-1} A_g^{\op})
\\
&= (g(\tilde \lambda) \frac{\xi^2}{g(\det(D_b))} A_g^{-1} B') \circ \Sigma \circ (B^{-1}A_g^{\op}) = A_g^{-1} \circ (\tilde \Sigma) \circ A_g^{\op}.
\end{align*}
By Lemma \ref{3-pointNonEq} and the fact that since $\Sigma = (yz,xz, xy)$, we have $\Sigma \circ \lambda I_3 = \lambda^2I_3 \circ \Sigma$, $\lambda \in \FF$, and, for any diagonal matrix $D$, we have $D \circ \Sigma \circ D = \det(D) I_3 \circ \Sigma$. This shows that $\tilde \Sigma$ is an affine representative of a $3$\nobreakdash-link from $S^{\op}$ to $S$.
\end{proof}

If we are looking for $[A_3] \sigma [A_2]$ that sends the base points of $\tilde \sigma^{-1}$ to the ones of~$\tilde \sigma$, we can search for $[B^{-1} A_3] \sigma [A_2 B']$ that fixes $\{ [1\sep 0\sep 0], [0\sep 1\sep 0], [0\sep 0\sep 1]\}$. This leads us~to\vspace*{-3pt}
\[
A_3 := B D_h ((A_2B')^{\Sigma})^{-1},
\]
where $D_h$ is a diagonal. We~want to find $D_h$ such that $A_3$ is the affine representative of an element of $\Aut_{\kk}(S^{\op})$.

\begin{lemma} \label{A3}
Let $\LL_p, \LL_q,\FF$ be as in our Convention \ref{con:2SF}, where $\kk \subset \LL_p, \LL_q$ are of degree $3$. There exists a $b \in \FF$ such that $N_{\FF/\LL_q}(b) = \xi^2$ and $B:= B_{S^{\op}}(b), B':= B_S(b^{-1})$.

Let $N$ be the affine representative of a $\kk$-automorphism of $S$. We~find a $\lambda_b \in \FF$ such that for $D_h = \diag(\lambda_b, g'(\lambda_b), g'(g'(\lambda_b)))$ (where $g'$ is as in Convention \ref{con:2SF})
\[
M := B D_h ((NB')^{\Sigma})^{-1}
\]
is an affine representative of a $\kk$-automorphism of $S^{\op}$.
\end{lemma}

\begin{proof}
The existence of $b$ follows from Lemma \ref{3-pointNonEq}. Note that $N_{\FF/\LL_q}(\sfrac{g(b)^3}{\xi^2}) = 1$ and using Hilbert's Theorem 90 we may choose $\lambda_b \in \FF$ such that\vspace*{-3pt}
\[
\frac{\lambda_b}{g(\lambda_b)} = \frac{g(b)^3}{\xi^2}.
\]
We define $D_h := \diag(\lambda_b, g'(\lambda_b),g'(g'(\lambda_b)))$ and find:
\begin{align*}
M^{g'} = B (123) D_h^{g'} ((N B' (123))^{\Sigma})^{-1} &= B (123) D_h^{g'} (123)^{-1} ((NB')^{\Sigma})^{-1}
\\
&= B D_h ((NB')^{\Sigma})^{-1} = M.
\end{align*}
This again shows that $M$ is defined over $\LL_p$. For the representation property, we~consider\vspace*{-3pt}
\begin{align*}
M &= \xi^2 A_g^{\op} B^g (D_b^{-1})^g D_h (D_{b^{-1}}^2)^g (((NB')^{\Sigma})^{-1})^g (A_g^{\op})^{-1}
\\
&= A_g^{\op} B^g D_h^g (((NB')^{\Sigma})^{-1})^g (A_g^{\op})^{-1} = A_g^{\op} M^g (A_g^{\op})^{-1},
\end{align*}
since\vspace*{-3pt}
\[
(D_b^{-1})^g D_h (D_{b^{-1}}^2)^g = \xi^{-2}D_h^g
\]
by our choice of $D_h$ and Lemma \ref{3-pointNonEq}. Therefore, $M$ is an affine representative of an automorphism in $\Aut_{\kk}(S^{\op})$.
\end{proof}

We next consider an example for the relation and use this to do our proof. We~use the construction of Lemma \ref{LargeCalc}, but this time some modifications are necessary since we work with $\tilde \sigma$ and not all scalars are defined over $\kk$.

\begin{lemma} \label{HardRel}
Let $\LL_p, \LL_q,\FF$ be as in our Convention \ref{con:2SF}, where $\kk \subset \LL_p, \LL_q$ are of degree $3$. Let $b \in \FF$ such that $N_{\FF/\LL_q}(b) = \xi^2$ and $B:= B_{S^{\op}}(b), B':= B_S(b^{-1})$. Let $\tilde \sigma$ be as in Lemma \ref{BB'}. We~define the following:
\begin{align*}
A_1 := I_3, \quad A_{2i} &:= D_{\xi^{-1}} (A_{2i-1}^{\tilde \Sigma})^{-1} = \tilde \lambda^{-1} D_{\xi^{-1}} M_2(B^{-1} A_{2i - 1})B'^{-1}, \quad i = 1,2
\\
A'_{6} &:= D_{\xi^{-1}} (A_{5}^{\tilde \Sigma})^{-1} = \tilde \lambda^{-1} D_{\xi^{-1}} M_2(B^{-1} A_{5})B'^{-1}
\\
\tag*{and}
A_{2i + 1} &:= B D_h ((A_{2i}B')^{\Sigma})^{-1} = B D_h M_2(A_{2i}B'), \quad i=1,2.
\end{align*}
Then for $\tau$ the $3$\nobreakdash-link represented by $\tilde \sigma$, $\alpha_i$ the automorphism represented by~$A_i$ and~$\alpha_6$ the automorphism represented by\vspace*{-3pt}
\[
A_6 := \tilde \lambda^{-3}D(B^{-1})^{-1}D_{\xi}^3 A'_6,
\]
we get the elementary relation\vspace*{-3pt}
\[
\alpha_6 \tau \dots \alpha_1 \chi = \id.
\]
\end{lemma}

\begin{proof}
We first prove that all the matrices are affine representatives. By Lemmas \ref{FixMatr} and \ref{A3}, the $A_i, A'_6$ are affine representatives. It remains to prove that $ \tilde \lambda^{-3} D(B^{-1})^{-1}D_{\xi}^3$ is also an affine representative.

We see that $D(B^{-1})$ is given by the coefficients of $B^{-1}$ as in Lemma \ref{LargeCalc}. We~have $(B^{-1})^{g'} = (123)^{-1}B^{-1}$, which is a permutation of the lines of $B^{-1}$. But using this, we~see $D(B^{-1})^{g'} = D(B^{-1})$. Therefore, using $\tilde \lambda \in \LL_p$, we~conclude that $ \tilde \lambda^{-3} D(B^{-1})^{-1}D_{\xi}^3$ is defined over $\LL_p$. Since\vspace*{-3pt}
\[
(B^{-1})^g = \xi (D_b^{-1})^g B^{-1}A_g^{\op}
\quand
D(A)_{ii} = \det(A)^6\,\frac{\prod_{k=1}^3(A^{-1})_{ik}^2}{\prod_{l=1}^3a_{li}}
\]
a direct calculation gives
\begin{align*}
D(B^{-1})^g_{11} &= D((B^{-1})^g)_{11} = D(\xi (D_b^{-1})^g B^{-1}A_g^{\op})_{11}
\\[-3pt]
&= D(B^{-1})_{22} (\xi^3 \det((D_b^{-1})^g)^3) = D(B^{-1})_{22} \frac{\xi^{-3} \tilde \lambda^3}{g(\tilde \lambda ^3)},
\\
D(B^{-1})^g_{22} &= D((B^{-1})^g)_{22} = D(\xi (D_b^{-1})^g B^{-1}A_g^{\op})_{22}
\\[-3pt]
&= D(B^{-1})_{33} (\xi^3 \det((D_b^{-1})^g)^3) =D(B^{-1})_{33} \frac{\xi^{-3} \tilde \lambda^3}{g(\tilde \lambda ^3)},
\\
D(B^{-1})^g_{33} &= D((B^{-1})^g)_{33} = D(\xi (D_b^{-1})^g B^{-1}A_g^{\op})_{33}
\\[-3pt]
&= D(B^{-1})_{11} (\xi^{12} \det((D_b^{-1})^g)^3) = D(B^{-1})_{11} \frac{\xi^{6} \tilde \lambda^3}{g(\tilde \lambda ^3)},
\end{align*}
since $(\sfrac{\tilde \lambda}{g(\tilde \lambda)})^3 = \sfrac{\xi^6}{\det(D_b^g)^3}$. Hence $\tilde \lambda^{3} D_{\xi}^{-3} D(B^{-1})$ is an affine representative.

We now show that $A_6 \tilde \sigma A_5 \sigma \dots A_1 \sigma = \id$. For this, we~relate $A_i$ to $M_i(B^{-1})$. To~simplify the calculation, we~will work with
\[
C_{2i} := A_{2i} B', \;i = 1,2,\quad C_{2i + 1} := B^{-1} A_{2i + 1}, \;i = 0,1,2, \quad C_{6} := A'_{6} B',
\]
for which we obtain
\begin{align*}
C_{2i} &= \tilde \lambda^{-1} D_{\xi^{-1}} M_2(C_{2i -1}), \quad i = 1,2,& C_{6} &= \tilde \lambda^{-1} D_{\xi^{-1}} M_2(C_{5})
\\
C_{2i + 1} &= D_h M_2(C_{2i}), \quad i=1,2, & C_1 &= B^{-1}.
\end{align*}
Thus we observe with Remark \ref{RemDiag}, that in $\Bir_{\FF}(\mathbb{P}^2)$:
\begin{align*}
A_6 \tilde \sigma A_5 \sigma \dots A_1 \sigma &= (D(B^{-1})^{-1}D_{\xi}^3) A'_6 B' \sigma B^{-1} A_5 \dots B^{-1}A_1 \sigma
\\
&= (D(B^{-1})^{-1}D_{\xi}^3) C_6\sigma \dots C_1 \sigma\\
&= (D(B^{-1})^{-1}D_{\xi}^3 D_{\xi^{-1}}) (D_{\xi^{-1}}^{-2}) (D_{\xi^{-1}}^4) D(B^{-1}) = \id.
\end{align*}
By Lemma \ref{ElemRel} and using $\varphi, \varphi^{\op}$, $\alpha_6 \tau \dots \alpha_1 \chi = \id$ is an elementary relation.
\end{proof}

Again, this example is enough to prove the case of Theorem \ref{MainA} where the splitting fields of the base points of $\chi$ and $\tau$ are both of degree $3$ over $\kk$.

\begin{theorem} \label{thm: elrel2}
Take an elementary relation between Sarkisov links of $S$, where both $3$\nobreakdash-point classes are non-equivalent and their splitting fields are degree-$3$ extensions of~$\kk$. Then the determinant of the elementary relation is trivial.
\end{theorem}

\begin{proof}
By Lemma \ref{DepPoint}, it suffices to look at all possible base points of $\tau \alpha_1$ and choose one version of the other $5$ automorphisms such that the base points are sent to each other by the links, as shown in Figure \href{https://jep.centre-mersenne.org/articles/10.5802/jep.347/}{1 in the original publication, p.1285}.

We choose $B,B', \tilde \sigma$ as in Lemma \ref{BB'}. The $3$\nobreakdash-points occurring in this relation cannot be on one of the three (geometric) lines between two of the three components of $p$ by the general position assumption we need for the elementary relation. Using the isomorphism $\varphi$ this means that the image of the $3$\nobreakdash-point does not intersect with the three lines $x,y,z = 0$. Therefore we can get all $3$\nobreakdash-points non equivalent to $p$ with $B$. We~therefore use the example relation from Lemma \ref{HardRel}.

We now calculate the matrices $A_i$ explicitly to be able to find the determinant, since in this case a priori we are only in $(\FF^*)^3$. We~calculate the matrices $A_i$ using the matrices $C_i$ from the proof of Lemma \ref{HardRel}. Note that for a diagonal matrix $D$, a general matrix $A$ and a scalar $\mu$ we get as in Remark \ref{RemDiag} that:
\[
M_2(A D) = D^{-2} M_2(A),
\quad
M_2(DA) = \det(D)^{-1} M_2(A) D,
\quad
M_2(\mu A) = \mu^{-2} M_2(A).
\]
Using this and $M_2 \circ M_2 = M_3$ we get (writing $M_i$ instead of $M_i(B^{-1})$):
\begin{align*}
C_2 &= \lambda_2 D_{\xi^{-1}} M_2, & \lambda_2 &= \tilde \lambda^{-1}
\\
C_3 &= \lambda_3 D_h M_3 D_{\xi^{-1}}, &\lambda_3 &= \tilde \lambda^2
\\
C_4 &= \lambda_4 D_{\xi} M_4 D_h, &\lambda_4 &= \tilde \lambda^{-5} \det(D_h^{-1})
\\
C_5 &= \lambda_5 D_h^{-1} M_5 D_{\xi}, &\lambda_5 &= \tilde \lambda^{10} \det(D_h)^2
\\
C_6 &= \lambda_6 D_{\xi^{-3}} M_6 D_h^{-1}, &\lambda_6 &=\tilde \lambda^{-22} \det(D_h)^{-3}.
\end{align*}
It remains to calculate the determinant and verify that it is actually contained in~$(\kk^*)^3$ and not only in $(\FF^*)^3$.

Using Lemma \ref{LargeCalc}, that:
\[
\det(\tilde \lambda^3 D(B^{-1})D_{\xi}^{-3}) = \tilde \lambda^9 \xi^{-9} \det(D(B^{-1})) = \tilde \lambda^9 P(B^{-1})^{-1} P(B)^2 \det(B^{-1})^{18}.
\]
Our goal is to show that
\[
\frac{\det(A_2A_4A'_6)\det(A_1A_3A_5)^{-1}} {\det(\tilde \lambda^{3}D(B^{-1}) D_{\xi}^{-3})} = \det(A_2A_4A_6)\det(A_1A_3A_5)^{-1}
\]
belongs to $(\kk^*)^3$. We~have
\begin{align*}
\det(A_2A_4A'_6)\det(A_1A_3A_5)^{-1} &= \det(C_2C_4C_6)\det(C_1C_3C_5)^{-1} \det(B')^{-3}\det(B)^3
\\
&= d(B^{-1}) \tilde \lambda^{-117} \det(D_h)^{-18}\det(B')^{-3}\det(B)^{-3}.
\end{align*}
Using Lemma \ref{LargeCalc} we obtain
\begin{multline*}
\frac{\det(\tilde \lambda^3 D(B^{-1})D_{\xi}^{-3})}{\det(A_2A_4A'_6)\det(A_1^{-1}A_3^{-1}A_5^{-1})} = \frac{\tilde \lambda^9 P(B^{-1})^{-1} P(B)^2 \det(B^{-1})^{18}}{d(B^{-1}) \tilde \lambda^{-117} \det(D_h)^{-18}\det(B')^{-3}\det(B)^{-3}}
\\
= P(B^{-1})^{-15} P(B)^9 \det(B^{-1})^{60} \det((B^{-1})^{\Sigma})^{42} \tilde \lambda^{126} \det(D_h)^{18}\det(B')^3\det(B)^3.
\end{multline*}
We now show that this is contained in $(\kk^*)^3$. To~this end, we~use the following properties:
\[
g(P(B^{-1})) = P(\xi (D_b^{-1})^gB^{-1}A_g^{\op}) = \xi^6 g(\det(D_b)^{-3}) P(B^{-1}) = (\sfrac{\tilde \lambda}{g(\tilde \lambda)})^3P(B^{-1}),
\]
Here, the first equality is by interchanging $P,g$ and using Lemma \ref{3-pointNonEq} on $B$, and the second equality is a property of $P$. Furthermore, $\tilde{\lambda} \in \LL_p$ and is therefore fixed by~$g'$. In~addition, we~have $(B^{-1})^{g'} = T B^{-1}$ for some permutation matrix $T$ by Lemma~\ref{3-pointNonEq}. This implies $g'(P(B^{-1})) = P(B^{-1})$. We~obtain
\begin{equation} \label{eq:1P}
P(B^{-1}) \tilde \lambda^3 \in \kk^*,
\end{equation}
and a similar calculation shows
\[
P(B)\tilde \lambda^{-3} \in \kk^*,
\]
which again implies
\begin{equation} \label{eq:2Ps}
P(B^{-1})P(B)\in \kk^*.
\end{equation}
We also see $g(\det(D_b))\xi^{-2} = \sfrac{g(\tilde \lambda)}{\tilde \lambda}$ by Lemma \ref{BB'} and obtain using Lemma \ref{3-pointNonEq}:
\[
g(\det(B)) = \xi^{-2} g(\det(D_b)) \det(B) = \frac{g(\tilde \lambda)}{\tilde \lambda}\, \det(B).
\]
Again we have $\tilde{\lambda}$ is fixed by $g'$. Using Lemma \ref{3-pointNonEq} we have $B^{g'} = B T$ for some permutation matrix $T$ of determinant $1$. This implies that $\det(B)$ is also fixed by $g'$, thus\vspace*{-3pt}
\begin{equation} \label{eq:1B}
\det(B) \tilde \lambda^{-1} \in \kk^*,
\end{equation}
and similarly\vspace*{-3pt}
\[
\det(B') \tilde \lambda \in \kk^*,
\]
and thus also\vspace*{-3pt}
\begin{equation} \label{eq:2Bs}
\det(B)\det(B') \in \kk^*.
\end{equation}
Since\vspace*{-3pt}
\[
\frac{\lambda_b}{g(\lambda_b)} = \frac{g(b)^3}{\xi^2}, \quad\frac{\tilde \lambda}{g(\tilde \lambda)} = \frac{\xi^2}{N_{\FF/\LL_p}(g(b))},
\]
it follows that\vspace*{-3pt}
\[
\frac{\det(D_h)}{g(\det(D_h))} = N_{\FF/\LL_p}(\sfrac{\lambda_b}{g(\lambda_b)}) = \frac{N_{\FF/\LL_p}(g(b))^3}{\xi^6} =\Bigl(\frac{g(\tilde \lambda)}{\tilde \lambda}\Bigr)^3.
\]
We thus get:\vspace*{-3pt}
\begin{equation} \label{eq:Dh}
\tilde \lambda^3 \det(D_h) \in \kk^*.
\end{equation}
Using these equations to reduce the determinant of the relation by elements of $(\kk^*)^3$. We~will make a sequence of equalities inside $\kk^*/(\kk^*)^3$ that all follow either from using an equation or reducing by $(\kk^*)^3$.

We have\vspace*{-3pt}
\begin{multline*}
\frac{\det(\tilde \lambda^3 D(B^{-1})D_{\xi}^{-3})}{\det(A_2A_4A'_6)\det(A_1^{-1}A_3^{-1}A_5^{-1})}
\\
= P(B^{-1})^{-15} P(B)^9 \det(B^{-1})^{60} \det((B^{-1})^{\Sigma})^{42} \tilde \lambda^{126} \det(D_h)^{18}\det(B')^3\det(B)^3.
\end{multline*}
From the fact that $A_2^{-1} = \tilde \lambda B' (B^{-1})^{\Sigma} D_{\xi}$, $\det(D_{\xi}) \in (\kk^*)^3$, $\tilde \lambda^{126} = (\tilde \lambda ^3)^{ 42}$ we obtain
\[= P(B^{-1})^{-15} P(B)^9 \det(A_2^{-1})^{42} \det(D_h)^{18}\det(B')^{-39}\det(B)^{-57}
\]
From \eqref{eq:2Ps}, \eqref{eq:2Bs}, $39=3 \cdot 13$, and $ \det(A_2^{-1})^{42} \in (\kk^*)^3$ we obtain:\vspace*{-3pt}
\[
= P(B^{-1})^{-24} \det(D_h)^{18}\det(B)^{-18}.
\]
From \eqref{eq:1P} we get, since $72 = 24 \cdot 3$,\vspace*{-3pt}
\[
= (\tilde \lambda)^{72} \det(D_h)^{18}\det(B)^{-18}.
\]
Using \eqref{eq:Dh} and $72 = 18 + 54 = 18 + 3 \cdot 3 \cdot 6$
\[
= (\tilde \lambda)^{18} \det(B)^{-18}.
\]
Using \eqref{eq:1B} we finally get the following:
\[
= 1.\qedhere
\]
\end{proof}

It remains to consider the case where one or both of the $3$\nobreakdash-points have a splitting field that is a degree-$6$ extension. We~can reduce this case to the previous one using mostly Galois theory.

\begin{theorem} \label{MainA}
Let $S$ be a non-trivial Severi--Brauer surface over $\kk$ and $S^{\op}$ its opposite Severi--Brauer surface. Let $\chi, \chi_i \colon S \dra S^{\op}$, $\tau, \tau_i \colon S^{\op} \dra S$ be Sarkisov links based on $3$\nobreakdash-points such that there exist $\alpha_i, \delta_i \in \Aut_{\kk}(S)$, $\gamma_i, \beta_i \in \Aut_{\kk}(S^{\op})$ such that
\[
\tau_i = \alpha_i \tau \beta_i, \quad\chi_i = \gamma_i \chi \delta_i,
\]
and the base points of $\chi, \tau$ are in general positions. Furthermore, if
\[
\tau_3 \chi_3 \tau_2 \chi_2 \tau_1 \chi_1 = 1,
\]
then we have
\[
\prod_{i=1}^3 \det(\alpha_i\delta_i)\det(\beta_i\gamma_i)^{-1} \in (\kk^*)^3.
\]
\end{theorem}

\begin{proof}
We take an elementary relation between Sarkisov links of $S$, with $p,q$ the two $3$\nobreakdash-points of the relation. Let $\LL_p, \LL_q$ be their splitting fields (we choose the same field if they are equivalent). Considering $([\LL_p\sep \kk], [\LL_q\sep \kk])$ we have proved the theorem in the $(3,3)$ case in theorems \ref{TrivialCase}, \ref{thm: elrel2}. In~the $(3,6)$ case, we~may invert the relation to switch $p,q$. Therefore, it remains to consider the $(6,3)$ and the $(6,6)$ cases.
\begin{enumerate}
\item We start with $(6,3)$. Since $\kk \subset \LL_p$ is a Galois extension of degree $6$, we~find a field $\LL_2$ such that $\kk \subset \LL_2 \subset \LL_p$, where the first extension is of degree $2$ and the second of degree $3$.

Assume $S_{\LL_2}$ would have $\LL_2$-points, then they would correspond to a $2$-point in~$S_{\kk}$, since the Galois group of this field extension is of size $2$. But this is impossible, thus~$S_{\LL_2}$ is a non-trivial Severi--Brauer surface over a perfect field.

After a base change to $\LL_2$ the point $q$ is still a closed point of degree $3$. It needs to be one orbit, since otherwise there would be closed points of order not divisible by $3$ on $S$ over $\LL_2$, a contradiction. We~just need to make sure that it still has a splitting field with field extension $3$. But $q$ splits over $\LL_2 \LL_q$ and since $[\LL_2\sep  \kk] = 2, [\LL_q\sep  \kk] = 3$ we have $\LL_2 \cap \LL_q = \kk$ which implies that $\LL_2 \subset \LL_2 \LL_q$ is of degree $3$. This implies that $\LL_2\LL_q$ is the splitting field of $q$ over $\LL_2$.

We now are in the $(3,3)$ case and can use the previous result. We~can use the isomorphism $\varphi, \varphi^{\op}$ for $\xi,\xi^{-1}$ from Lemma \ref{3pLem} that we have always used, since for degree-$6$ extensions they give the same $A_g$ and we may keep our affine representation of the automorphisms. Thus, by the theorems in the $(3,3)$ cases, we~conclude that the determinant of the word is in $\kk^* \cap (\LL_2^{*})^3$. But now it is only left to see that \hbox{$\kk^* \cap (\LL_2^{*})^3 \subset (\kk^{*})^3$}.

For this, choose $a \in \kk^*$ and $b \in \LL_2^*$ such that $a = b^3$. Then if the polynomial $x^3 - a$ has a solution in $\kk$ we are done. If not, we~see the minimal polynomial of $b$ over $\kk$ is of degree $3$, which is impossible for a degree-$2$ field extension. This finishes this case.
\item For $(6,6)$ we again take $\LL_2$ so that $\kk \subset \LL_2 \subset \LL_p$ with the same degrees as above. As above $S_{\LL_2}$ is still a non-trivial Severi--Brauer surface over a perfect field and $q$ is still a $3$\nobreakdash-point (its splitting field can be of both degree $3$ or $6$). We~now are either in the $(3,3)$ case or in the $(3,6)$ case.

In the $(3,3)$ case the result can be shown as above using the same $\varphi, \varphi^{\op}$ and $\kk^* \cap (\LL_2^{*})^3 \subset (\kk^{*})^3$.

In the $(3,6)$ case we again use the same $\varphi, \varphi^{\op}$ and see by the previous case that the determinant of the word is in $ (\LL_2^{*})^3$. However, because of the use of the same isomorphisms, we~see it is also an element of $\kk^*$. We~now argue as in the end of the $(6,3)$ case.\qedhere
\end{enumerate}
\end{proof}

We now have all the properties needed to extend the determinant homomorphism to our group of birational self-maps. We~can also extend it to the groupoid $G_S$ introduced in the beginning.

\begin{corollary}\label{DetOnBir}
For every choice $R$ of representatives of the different classes of $\kk$-links (from $S$ to $S^{\op}$) there exists a well defined non-trivial groupoid homomorphism
\begin{align*}
\det_R \colon G_S &\to \kk^*/(\kk^*)^3
\\
R &\mto 1
\\
\alpha \in \Aut_{\kk}(S) &\mto \det(\alpha)
\\
\beta \in \Aut_{\kk}(S^{\op}) &\mto \det(\beta)^{-1},
\end{align*}
where $G_S := \{ \varphi \colon S_1 \dra S_2 : S_1,S_2 \in \{ S, S^{\op}\}$, $\varphi \text{ birational}\}$.
\end{corollary}

\begin{proof}
From Theorem \ref{Isk} and Lemma \ref{Links36} we get that every $G_S$ is generated by links from $R$ and automorphisms of $S,S^{\op}$, with elementary relations given by Lemma~\ref{ElRel} and trivial relations given by
\[
\chi^{-1} \chi = \id, \quad\chi^{-1} \beta \chi \alpha = \id.
\]
For $\varphi \in G_S$ given by
\[
\varphi = \alpha_{r+1}\prod_{i=1}^r\chi_i \alpha_i,
\]
where $\chi_i \in R$ and $\alpha_i$ are automorphisms of $S$ if $i$ is odd and automorphisms of $S^{\op}$ if $i$ is even, we~define
\[
\det_R(\varphi) := \det(\alpha_1) \det(\alpha_2^{-1}) \dots
\]
The first trivial relation is sent to the identity by definition, the second trivial relation is sent to the identity by Lemma \ref{DetTrivRef}. The elementary relation is sent to the identity by Theorem \ref{MainA}. To~show that it is non-trivial we take a $3$\nobreakdash-point $p$ whose splitting field $\LL_p$ is a degree-$3$ extension over $\kk$ (exists, as mentioned above) and hence, using the usual isomorphism $\varphi$ we obtain the corresponding element $\xi$ which is not a norm and hence does not lie in $(\kk^*)^3$. This $\xi$ is the determinant of $A_g$, and it remains to prove that $A := A_g$ is the affine representative of some automorphism of $S$.

The matrix $A$ is an affine representative of an automorphism of $S$ if and only if
\[
A^g = A_g^{-1} A A_g,
\]
which, as in Lemma \ref{AutRep} comes from $A$ commuting with $A_g \circ g$ which is the formal definition of $A$ descending to an automorphism of $S$. But the above equation is true since $A_g^g = A_g = A_g^{-1}A_g A_g$.
\end{proof}

\begin{thebibliography}{99}
\bibitem[Ser97]{serre1997galois} J.-P.Serre, \emph{Galois cohomology}, Springer Monographs in Math., Springer-Verlag,Berlin,1997.
\bibitem[Ser04]{serre2004corps} J.-P.Serre, \emph{Corps locaux}, Actualités scientifiques et industrielles, Hermann,Paris,2004.
\bibitem[BSY24]{BSY} J.Blanc,J.Schneider and E.Yasinsky, ``Birational maps of Severi-Brauer surfaces,with applications to Cremona groups of higher rank'',2024, arXiv:2211.17123.
\bibitem[BSY]{BSYUnp} J.Blanc,J.Schneider and E.Yasinsky, ``Unpublished notes''. Original author's attribution; the complete licensed published proofs used here are retained.
\bibitem[LZ20]{Lamy_2020} S.Lamy and S.Zimmermann, ``Signature morphisms from the Cremona group over a non-closed field'', \emph{J. Eur. Math. Soc.}(JEMS)\textbf{22}(2020),no.10,3133--3173.
\bibitem[Isk96]{Iskovskikh} V.A.Iskovskikh, ``Factorization of birational maps of rational surfaces from the viewpoint of Mori theory'', \emph{Russian Math. Surveys}\textbf{51}(1996),no.4,585--652.
\bibitem[Shr20]{zbMATH07206468} C.A.Shramov, ``Birational automorphisms of Severi-Brauer surfaces'', \emph{Mat. Sb.}\textbf{211}(2020),no.3,466--480.
\bibitem[Wei22]{cm:9040} F.Weinstein, ``On birational automorphisms of Severi-Brauer surfaces'', \emph{Commun. Math.}\textbf{30}(2022),no.1,article1(9pages).
\bibitem[Kol25]{kollár2025severibrauervarietiesgeometrictreatment} J.Kollár, ``Severi-Brauer varieties; a geometric treatment'',2025, arXiv:1606.04368.
\bibitem[Jah03]{jahnel} J.Jahnel, \emph{The Brauer-Severi variety associated with a central simple algebra: a survey}, Math. Gottingensis2, Mathematisches Institut, Universität Göttingen,Göttingen,2003.
\bibitem[AC02]{zbMATH01683518} M.Alberich-Carramiñana, \emph{Geometry of the plane Cremona maps}, Lect. Notes in Math.1769, Springer,Berlin,2002.
\end{thebibliography}
\end{document}
